% concatenated from Symplectic_self_orthogonal_quasi_cyclic_codes_final_version.tex
%\documentclass[journal,onecolumn]{IEEEtran}
\documentclass[journal]{IEEEtran}
%\documentclass[journal,onecolumn]{IEEEtran}
%\usepackage{geometry}
%\geometry{a4paper,left=2.5cm,right=2.5cm,top=3cm,bottom=3cm}
%\linespread{2}
%\documentclass[journal,onecolumn]{IEEEtran}
\usepackage{amsmath,amsfonts}
\usepackage{algorithmic}
\usepackage{algorithm}
\usepackage{array}
\usepackage[caption=false,font=normalsize,labelfont=sf,textfont=sf]{subfig}
\usepackage{textcomp}
\usepackage{stfloats}
\usepackage{url}
\usepackage{verbatim}
\usepackage{graphicx}
\usepackage{amsmath}
\usepackage{amssymb}
\usepackage{cuted} %公式单栏排双栏
\usepackage{mathrsfs}
\usepackage{makecell}
\usepackage{threeparttable}
\usepackage[figuresright]{rotating}
\usepackage{arydshln}
\usepackage{setspace}%使用间距宏包
\usepackage{color, xcolor} % 颜色包，color 必须导入，xcolor 建议导入
\usepackage{cite}
%\usepackage[UTF8]{ctex}
\newtheorem{lemma}{Lemma}
\newtheorem{theorem}{Theorem}
\newtheorem{corollary}{Corollary}
\newtheorem{definition}{Definition}
\newenvironment{proof}{{\noindent\it Proof:} }{\hfill $\square$\par}
\newtheorem{proposition}{Proposition}
\newtheorem{example}{Example}
\newtheorem{remark}{Remark}
\newenvironment{note}{{\it Note:}}{}
\hyphenation{op-tical net-works semi-conduc-tor IEEE-Xplore}
% updated with editorial comments 8/9/2021



\begin{document}

\title{Symplectic self-orthogonal quasi-cyclic codes}

\author{Chaofeng Guan, Ruihu Li, Jingjie Lv, Zhi Ma~\IEEEmembership{}
%\thanks{Chaofeng Guan is with Henan Key Laboratory of Network Cryptography Technology, Zhengzhou, 450001, China.  (e-mail: gcd2020@yeah.net).}
%\thanks{Ruihu Li is with the Fundamentals Department, Air Force Engineering University, Xi’an, 710051, China. (e-mail: liruihu@aliyun.com).}
%\thanks{Jingjie Lv is with the 
%	School of Electrical Engineering \& Intelligentization, Dongguan University of Technology, Dongguan, 523808,
%	China (email: juxianljj@163.com).}
%\thanks{Zhi Ma is with Henan Key Laboratory of Network Cryptography Technology, Zhengzhou, 450001, China (e-mail: LWtglx2023@outlook.com).}
%\thanks{Chaofeng Guan is supported by the National Key Research and Development Program of China under Grant 2021YFB3100100. 
%	Ruihu Li is supported by the National Natural Science Foundation of China under Grant U21A20428.
%}
}
%61972413,61901525,62002385,U21A20428
\markboth{}%
{}

\IEEEpubid{}
% Remember, if you use this you must call \IEEEpubidadjcol in the second
% column for its text to clear the IEEEpubid mark.

\maketitle

\begin{abstract}
%we propose a method for constructing symplectic self-orthogonal quasi-cyclic codes, which allows arbitrary polynomials that coprime $x^{n}-1$ to construct symplectic self-orthogonal codes. Moreover, 
In this paper, we establish the necessary and sufficient conditions for quasi-cyclic (QC) codes with index even to be symplectic self-orthogonal. Subsequently, we present the 
lower and upper bounds on the minimum symplectic distances of a class of $1$-generator QC codes and their symplectic dual codes by decomposing code spaces.
As an application, we construct numerous new binary symplectic self-orthogonal QC codes with excellent parameters, leading to $117$ record-breaking quantum error-correction codes. 
%improving the records in Grassl's online table (Bounds on the minimum distance of QECCs, http://www.codetables.de). 
\end{abstract}

\begin{IEEEkeywords}
Quasi-cyclic code, symplectic self-orthogonal code, record-breaking quantum code.
\end{IEEEkeywords}


\section{Introduction}
\IEEEPARstart{Q}{uantum} error-correction codes (QECCs) are capable of encoding quantum bits to overcome unwanted noise and decoherence and are considered one of the cornerstones of quantum computing and quantum communication. 
The concept of QECCs was initially proposed by Shor \cite{shors1995scheme} and Steane \cite{steane1996multiple}.
Subsequently, many scholars have studied QECCs and gradually established the connection between QECCs and classical codes with self-orthogonal properties under Euclidean, Hermitian, and symplectic inner products \cite{calderbank1998quantum,ketkar2006nonbinary,steane1999enlargement}.
Later, a large number of QECCs with suitable parameters were constructed using classical codes such as cyclic codes \cite{aly2007quantum,shi2019dual}, constacyclic codes \cite{kai2014constacyclic,chen2015application}, and algebraic geometry codes \cite{jin2011euclidean,bartoli2021certain}.
However, due to the inherent limitations of these classical codes, the constructed QECCs are often limited to specific lengths and dimensions. Therefore, these good QECCs are discrete.
Specifically, for binary QECCs, improving the existing records in Grassl's online code table \cite{Grassltable} faces two significant challenges. 
First, many researchers have tried exhaustive searches when the lengths are small. Without additional clever strategies, it is unlikely to find better codes. Second, for QECCs with long lengths, the computational complexity of determining their actual distances will be very high.
Therefore, more powerful tools are needed to construct QECCs with excellent parameters. 


As an effective generalization of cyclic codes, quasi-cyclic (QC) codes possess a rich algebraic structure that allows for the construction of flexible codes in terms of lengths and dimensions. In 1974, Kassmi \cite{kasami1974gilbert} proved that QC codes satisfy the modified Gilbert-Varshamov bound, thus indicating their good asymptotic performance. In recent decades, many scholars have conducted in-depth research on the algebraic structure of QC codes \cite{ling2001algebraic,guneri2012concatenated,semenov2012spectral,ezerman2021comparison}, 
resulting in the development of numerous record-breaking linear codes from QC codes \cite{aydin2001structure,daskalov2003new,Chen2016NewBH,aydin2019equivalence}. These studies have demonstrated the excellent performances of QC codes. Therefore, it is natural to consider using QC codes to construct high-quality QECCs. 
%As an effective generalization of cyclic codes, the abundant algebraic structure of quasi-cyclic codes allows the construction of flexible codes in terms of lengths and dimensions.
%In 1974, Kassmi \cite{kasami1974gilbert} proved that quasi-cyclic codes satisfy the modified Gilbert-Vashamov (GV) bound. 
%Therefore, the parameters of quasi-cyclic codes have asymptotic good performance.
%In recent decades, many scholars have studied the algebraic structure of quasi-cyclic codes in depth \cite{ling2001algebraic,guneri2012concatenated,semenov2012spectral,ezerman2021comparison} and constructed a large number of record-breaking linear codes from quasi-cyclic codes \cite{aydin2001structure,daskalov2003new,Chen2016NewBH,aydin2019equivalence}, which demonstrated their excellent performance. Therefore, it is natural to consider using quasi-cyclic codes to construct good QECCs. 
%Hagiwara et al. \left[12,13\right] studied constructions of quantum codes by QC LDPC
%codes. They focused on long codes and probabilistic constructions.
Work on constructing QECCs from QC codes started with Hagiwara et al. \cite{hagiwara2007quantum}, but progress was slow due to the lack of a practical theory. 
In 2018, Galindo et al. \cite{galindo2018quasi} determined some sufficient conditions for dual-containment of a class of QC codes under Euclidean, Hermitian, and symplectic inner products. 
Although they did not construct record-breaking QECCs, their work provided a new perspective.
Inspired by Galindo's work, Ezerman et al. \cite{ezerman2019good} studied QC codes over finite fields of order $4$ and $9$ and obtained some new binary and ternary QECCs.
Lv et al. \cite{lv2020explicit,lv2020constructions} and the authors \cite{guan2022construction,guan2022new} studied sufficient conditions for dual-containing or self-orthogonal of three classes of QC codes with specific algebraic structures, and constructed many record-breaking binary QECCs. These studies demonstrate that QC codes also possess superior properties in constructing QECCs.





Nevertheless, many important issues remain concerning constructing QECCs from QC codes. The crucial issue is identifying what kind of QC codes are self-orthogonal, so one can construct them efficiently. 
Notably, this paper determines, for the first time, the necessary and sufficient conditions for symplectic self-orthogonality of QC codes with index even.
The second problem is determining the bounds on the minimum distance of QC codes. 
Around this problem, this paper proposes lower and upper bounds on the minimum symplectic distance of a class of QC codes with index 2, which improves the results in \cite{dastbasteh2023polynomial}, \cite{GuanAdditive2023} and \cite{galindo2018quasi}. 
As a practical application of our findings, we construct numerous binary QECCs that break previous records in terms of their error-correcting capabilities.
%As an application, we construct many record-breaking binary QECCs. 
%Notably, in the case of $1$-generator quasi-cyclic codes with index 2, our symplectic self-orthogonal  codes construction is different from Calderbank's additive construction (Theorem 14 in \cite{calderbank1998quantum}).



% In Section 3, we discuss the necessary and sufficient conditions for quasi-cyclic codes of index even  to be symplectic self-orthogonal, and provide practical approaches for constructing quasi-cyclic symplectic self-orthogonal codes. In Section 4, we propose a class of symplectic self-orthogonal quasi-cyclic codes and construct many record-breaking binary QECCs. Finally, in Section 5, we discuss our main results and future research directions.
This paper is structured as follows. Section \ref{sec2} provides some basics that we use in this paper.
In Section \ref{sec3}, we discuss the necessary and sufficient conditions for QC codes of index even to be symplectic self-orthogonal.
In Section \ref{sec4}, we provide bounds on the minimum symplectic distances of a class of $1$-generator QC codes and their symplectic dual codes. Then, we give some interesting and record-breaking results.  
Finally, we conclude this paper in Section \ref{sec5}. 
%This paper presents the parameters of symplectic self-orthogonal or dual-containing codes as $\left[2n,k\right]_q^s$ or $\left[2n,k,d_s\right]_q^s$, where $d_s$ denotes symplectic distance.
\section{Preliminaries}\label{sec2}
This section presents the basics of classical codes and QECCs to facilitate the development of subsequent sections. For more details, see standard textbooks 
\cite{huffman2010fundamentals,huffman2021concise}.
%This section presents the fundamentals of classical codes and QECCs to facilitate the development of subsequent sections. For more detailed information, we refer the reader to standard textbooks \cite{huffman2010fundamentals,huffman2021concise}.

\subsection{Linear codes}

Let $\mathbb{F}_q$ be a finite field with $q$ elements, where $q$ is a prime power. Let $\mathbb{F}_q^*=\mathbb{F}_q\setminus \{0\}$. % and positive integer $\ell$ satisfies $\ell \equiv 0\pmb{2}$.
An $\left[\ell n,k\right]_q$ linear code $\mathscr{C}$ is a $k$-dimensional subspace of  $\mathbb{F}_{q}^{\ell n}$. 
For a codeword $\mathbf{c}=(c_{0},c_1,\ldots, c_{\ell n-1})$ of $\mathscr{C}$, the Hamming weight of $\mathbf{c}$ is $\mathbf{w}_{H}(\mathbf{c})=|\left\{i \mid c_{i} \neq 0, 0 \leq i \leq \ell n -1 \right\}|$.
If $\ell n$ is even and $N=\ell n/2$,
then the symplectic weight of $\mathbf{c}$ is 
$\mathbf{w}_{s}(\mathbf{c})=|\{i \mid (c_{i}, c_{N+i})  \neq(0,0), 0 \leq i \leq N-1 \}|$. 
The support of $\mathbf{c}$, denoted by $\mathrm{Supp}(\mathbf{c})$, is the set of coordinates of nonzero entries in $\mathbf{c}$, which is a subset of $\{0,1,\ldots,\ell n-1\}$.
If the minimum Hamming (resp. symplectic) weight of $\mathscr{C}$ is $d_{H}=\min \left\{\mathbf{w}_{H}(\mathbf{c}) \mid \mathbf{c} \in \mathscr{C}, \mathbf{c} \ne \mathbf{0} \right\}$ (resp. $d_{s}=\min \left\{\mathbf{w}_{s}(\mathbf{c}) \mid \mathbf{c} \in \mathscr{C}, \mathbf{c} \ne \mathbf{0}\right\}$), then $\mathscr{C}$ can be written as an $\left[\ell n,k,d_H\right]_q$ (resp. $\left[\ell n,k,d_s\right]_q^s$ ) code.


For $\mathbf{u }=(u_{0},u_1,  \ldots, u_{\ell n-1})$, 
$\mathbf{v}=(v_{0},v_1,\ldots, v_{\ell n-1}) \in \mathbb{F}_q^{\ell n}$,  the Euclidean and symplectic inner products of them are  $\left\langle\mathbf{u}, \mathbf{v} \right\rangle_{e}=\sum\limits_{i=0}^{\ell n-1} u_{i} v_{i}$ and 
$\left\langle\mathbf{u},\mathbf{v} \right\rangle_{s}=\sum\limits_{i=0}^{N-1}\left(u_{i} v_{N+i}-u_{N+i} v_{i}\right)$, respectively. 
The Euclidean and symplectic dual codes of $\mathscr{C}$ are  $\mathscr{C}^{\perp_{e}}=\left\{\mathbf{v} \in \mathbb{F}_{q}^{\ell n} \mid\left\langle\mathbf{u}, \mathbf{v} \right\rangle_{e}=0, \forall \mathbf{u} \in \mathscr{C}\right\}$ and
$ \mathscr{C}^{\perp_{s}}= \left\{\mathbf{v} \in \mathbb{F}_{q}^{\ell n} \mid\left\langle\mathbf{u},\mathbf{v} \right\rangle_{s}=0, \forall \mathbf{u} \in \mathscr{C}\right\}.$
If $\mathscr{C} \subset \mathscr{C}^{\perp_{e}}$ (resp. $\mathscr{C} \subset \mathscr{C}^{\perp_{s}}$), then $\mathscr{C}$ is a Euclidean (resp. symplectic) self-orthogonal code. 


For $i \in\{1, 2\}$, let $\mathscr{C}_i$ be an $\left[n_i , k_i , d_i \right]_q$ code. Then, the \textit{ direct sum}  of $\mathscr{C}_1$ and $\mathscr{C}_2$ is an $\left[n_1 + n_2, k_1 + k_2, \min\{d_1, d_2\}\right]_q$ code, denoted by $\mathscr{C}_1\oplus  \mathscr{C}_2$.
The following lemma helps identify symplectic self-orthogonal codes.

\begin{lemma}\label{lemma-self-orthogonal}(\cite{huffman2021concise})
	Let $\ell$ be an even integer and $\mathscr{C}$ be a linear code over $\mathbb{F}_q^{\ell n}$. Then, $\mathscr{C}$ is a self-orthogonal code under the symplectic inner product if and only if for all $\mathbf{c_1}, \mathbf{c_2} \in \mathscr{C}$, it holds that $\left\langle \mathbf{c_1}, \mathbf{c_2}\right\rangle_s = 0$.
\end{lemma}
%\subsection{Additive codes}
%A code $\mathscr{C}_a$ is said to be an additive code over $\mathbb{F}_{q^2}$ if it is a subgroup of $\mathbb{F}^n_{q^2}$, which means that scalar multiples of the codewords do not necessarily belong to the code. 
%If $\mathscr{C}_a$ has dimension $\frac{k_a}{2}$, the minimum Hamming distance $d$, then $\mathscr{C}_a$ can be denoted as $(n,\frac{k_a}{2},d)_{q^2}$.
%Define a map $\Phi$ form $\mathbb{F}_q^{2n}$ to $\mathbb{F}_{q^2}^n$ as $\Phi(\mathbf{v})=(v_0+w\cdot v_n,v_1+w\cdot v_{n+1},\cdots,v_{n-1}+w\cdot v_{2n-1})$, where $w$ generates $\mathbb{F}_{q^2}$ over $\mathbb{F}_{q}$.
%Then, it is easy to see that $\Phi$ is a bijection, and $\mathbf{w}_H(\Phi(\mathbf{v}))=\mathbf{w}_{s}(\mathbf{v})$.
%Let $G$ denote generator matrix of $\mathscr{C}$ with parameters $\left[\ell n,k,d_s\right]_q^s$, then $\Phi(G)$ can generate additive code $\mathscr{C}_a$ with parameters $(\ell n,\frac{k}{2},d_s)_{q^2}$. In addition, By \cite{calderbank1998quantum}, a trace Hermitian self-orthogonal additive code over $\mathbb{F}_4$ is equivalent to a symplectic self-orthogonal code over $\mathbb{F}_2$. 
%Since additive code is not the focus of this paper, we will not go into it here. For more information, we refer the readers to Ref. .  
 

\subsection{Cyclic codes and quasi-cyclic codes}

Let $\left\langle x^{n}-1\right\rangle$ represent the ideal of the polynomial ring $\mathbb{F}_q\left[x\right]$ generated by $x^n -1$ and $\mathbb{R}_{q,n} = \mathbb{F}_q\left[x\right]/\left\langle x^{n}-1\right\rangle$ be the quotient ring of $\mathbb{F}_q\left[x\right]$ modulo $\left\langle x^{n}-1\right\rangle$.
An element in  $\mathbb{R}_{q,n}$ is an equivalence class of polynomials in $\mathbb{F}_q[x]$. Each equivalence class has a representative which is a polynomial in $\mathbb{F}_q[x]$ of degree less than $n$. Here, equivalence classes in $\mathbb{R}_{q,n}$ are represented by this representative. For a polynomial $c(x)$ in $\mathbb{F}_q[x]$, we use $[c(x)]$ to denote its representative.



A linear $\left[n,k\right]_q$ code $\mathscr{C}$ is cyclic if and only if for any codeword $\mathbf{c}=(c_{0}, c_{1}, \ldots, c_{n-1})$ of $\mathscr{C}$, the shifted vector $\tau_1(\mathbf{c})=(c_{n-1}, c_{0}, \ldots, c_{n-2})$ is also a codeword of $\mathscr{C}$. 
Define an isomorphic mapping $\pi$ as follows:  
\begin{equation}  
\pi: \mathbb{R}_{q,n} \rightarrow \mathbb{F}_q^n, \quad [c(x)]=\sum_{i=0}^{n-1} c_{i} x^{i} \mapsto \mathbf{c}=(c_0,c_1,\ldots,c_{n-1}).  
\end{equation}

Then, we have $\pi(xc(x))=\tau_1(\mathbf{c})$. 
%For $c(x)$ in $\mathbb{R}_{q,n}$, $\left[c(x)\right]$ denotes its equivalence class in  $\mathbb{R}_{q,n}$. 
%$\pi^{-1}(\tau_1(\mathbf{c}))$ corresponds to the class $\left[xc(x)\right]$. 
%We default that the representative of an equivalence class is the polynomial of least degree in its class. 
Therefore, when studying cyclic codes, $\left[c(x)\right]$ is identified with the coefficient vector $\mathbf{c}$.  
%Consider each codeword $\mathbf{c}$ as coefficient vector of a polynomial $\mathbf{c}(x)=\sum\limits_{i=0}^{n-1} c_{i} x^{i}$ in $\mathbb{F}_{q}\left[x\right]$. 
Since a cyclic code is an ideal in the quotient ring $\mathbb{R}_{q,n}$ and $\mathbb{R}_{q,n}$ is principal, every cyclic code can be generated by a single polynomial in $\mathbb{R}_{q,n}$.
For an $\left[n,k\right]_q$ cyclic code $\mathscr{C}$, we call the monic polynomial $g(x)$ with least degree $n-k$ that generates $\mathscr{C}$ as the generator polynomial of $\mathscr{C}$. 
We use $d_H(\mathscr{C})$ or $d(\left[g(x)\right])$ to denote the minimum Hamming distance of $\mathscr{C}$. 
The parity check polynomial of $\mathscr{C}$ is $h(x)=(x^{n}-1)/g(x)$, whose reciprocal polynomial is \( \tilde{h}(x)= x^{\deg(h(x))} h(x^{-1})\).
The Euclidean dual of $\mathscr{C}$ is also a cyclic code, which has generator polynomial $g^{\perp_e}(x)=h_0^{-1} \tilde{h}(x)$, where $h_0$ is the constant term coefficient of $h(x)$. If $g^{\perp_e}(x)\mid g(x)$, then $\mathscr{C}$ is a Euclidean self-orthogonal code.



% Note that $\mathbf{c}^{\prime}=(c_{\ell n-\ell}, c_{\ell n-\ell+1}, \ldots, c_{\ell n-1}, c_0, c_1, \ldots, c_{\ell n-\ell-1})$ is an alternative but equivalent way of writing $\tau_{\ell}(\mathbf{c})$.  
% Also, if $\ell=1$, then $\mathscr{C}$ is also a cyclic code.  
  


A linear code $\mathscr{C}$ of length $\ell n$ over $\mathbb{F}_{q}$ is called a QC code of index $\ell$, if for any codeword $\mathbf{c}=\left(c_{0}, c_{1}, \ldots, c_{\ell n-1}\right)$ of $\mathscr{C}$, the $\ell$ shifted vector  
\[  
\tau_{\ell}(\mathbf{c})=(c_{\ell n-\ell}, \ldots, c_{\ell n-1}, c_0, \ldots, c_{\ell n-\ell-1})  
\]  
is also a codeword.  
%Note that if $\ell=1$, $\mathscr{C}$ is also a cyclic code.
By \cite{huffman2021concise}, after appropriate permutations, the generator matrix $G$ of an $h$-generator QC code with index $\ell$ can be put into the following form:
\begin{equation}\setlength{\arraycolsep}{1.2pt}
	G=\left(\begin{array}{cccc}
		A_{1,0} & A_{1,1} & \ldots & A_{1, \ell-1} \\
		A_{2,0} & A_{2,1} & \ldots & A_{2, \ell-1} \\
		\vdots & \vdots & \ddots & \vdots \\
		A_{h, 0} & A_{h, 1} & \ldots & A_{h, \ell-1}
	\end{array}\right),
\end{equation}
where  $A_{i, j}$ $(1\le i \le h, 0\le j \le \ell-1)$ is an $n\times n$ circulant matrix. % established by polynomial $[a_{i,j}(x)] \in \mathbb{R}_{q,n}$. 
Therefore, a QC code of index $\ell$ also can be identified with an $\mathbb{R}_{q,n}$-submodule of $\mathbb{R}_{q,n}^{\ell}$.

%For $\ell \ge 2$ and $\ell$ is, quasi-cyclic codes also can be transformed into additive cyclic codes over $\mathbb{F}_{q^{2}}$ via map $\Phi$ above. 
%Therefore, quasi-cyclic codes in Section \ref{sec4} also can be viewed as additive cyclic codes over $\mathbb{F}_{4}$.
% generated by polynomials  $a_{i, j}(x)\in \mathbb{R}_{q,n}$, where  $1 \leq i \leq h$  and  $0 \leq j \leq \ell-1$. 
%
%
%Circulant matrices are basic components in the generator matrix for quasi-cyclic codes. An $n\times n$ circulant matrix $A$ is defined as
%\begin{equation}
%	A=\left(\begin{array}{ccccc}
%		a_{0} & a_{1} & a_{2} & \ldots & a_{n-1} \\
%		a_{n-1} & a_{0} & a_{1} & \ldots & a_{n-2} \\
%		\vdots & \vdots & \vdots & \vdots & \vdots \\
%		a_{1} & a_{2} & a_{3} & \ldots & a_{0}
%	\end{array}\right).
%\end{equation}
%
%If the first row of  $A$  is mapped onto polynomial  $a(x)$, then circulant matrix  $A$  is isomorphic to polynomial $a(x)=a_{0}+a_{1} x+\cdots+a_{n-1} x^{n-1} \in \mathbb{R}_{q,n}$. So $A$ can be determined by polynomial $a(x)$. 
%Generator matrix of $h$-generator quasi-cyclic code with index $\ell$ 
%%can be transformed into rows of $n\times n$ circulant matrices by suitable permutation of columns, which 
%has the following form:



\subsection{Quantum error-correction codes}
Since our purpose is to improve the records in \cite{Grassltable}, here we only introduce the basics of binary QECCs. 
Similar to the classical situation, a binary QECC $\mathrm{Q}$ of length $n$ is a $K$-dimensional subspace of $2^n$-dimensional Hilbert space $(\mathbb{C}^{2})^{\otimes n}$, where $\mathbb{C}$ represents complex field and $(\mathbb{C}^{2})^{\otimes n}$ is the $n$-fold tensor power of $\mathbb{C}^2$. If $K = 2^k$, then the parameters of $\mathrm{Q}$ are denoted by $\left[\left[n, k, d\right]\right]_2$, where $d$ is its minimum distance, and $\mathrm{Q}$ can correct up to $\left\lfloor\frac{d-1}{2}\right\rfloor$ qubit errors. 
\begin{lemma}\label{CRSS}(\cite{calderbank1998quantum}) If $\mathscr{C} \subset \mathbb{F}_{2}^{2 n}$ is a symplectic self-orthogonal $\left[2n, k\right]^s_2$ code, then there exists a QECC with parameters $\left[\left[n, n-k, d_s\right]\right]_{2}$, where $d_s$ is the minimum symplectic distance of $\mathscr{C}^{\perp_{s}} \backslash\mathscr{C}$.
\end{lemma}

In general, QECCs can also be derived from existing ones by the following propagation rules, which will help us to get more record-breaking QECCs later.

\begin{lemma}\label{propagation_rules}
	(\cite{calderbank1998quantum}) Suppose there exists an $\left[\left[n, k, d\right]\right]_{2}$ QECC. Then, the following QECCs also exist.\\
	(1) $\left[\left[n, k-1, d\right]\right]_{2}$ for $k\ge 1$;\\
	(2) $\left[\left[n+1, k, d\right]\right]_{2}$ for $k>0$;\\
	(3) $\left[\left[n-1, k+1, d-1\right]\right]_{2}$ for $n\ge2$.
\end{lemma}

\begin{note}
Throughout the paper, we agree on the following notation.
All the indices $\ell$ of QC codes in the subsequent texts will be even and set $m=\ell/ 2$.
For $f(x)=f_{0}+f_{1} x+\cdots+ f_{n-1} x^{n-1} \in \mathbb{R}_{q,n}$, denote $\bar{f}(x)=x^{n}f(x^{-1})$. 
%In this work, all polynomials are computed over $\mathbb{R}_{q,n}$, therefore we agree with that 
% \textcolor{blue}{When we consider the codeword or the inner product between polynomials, $\left[f(x)\right]$ is identified with the coefficient vector $\pi(f(x))=({f_{0}},f_{1},\ldots ,f_{n-1})$. }
Let $\mathbf{0}$ and $\mathbf{0}_n$ denote a appropriate zero matrix and a zero vector of length $n$, respectively.
Additionally, we denote 
$\Omega _{mn}=\left(\begin{array}{cc}
	\mathbf{0} & I_{m n} \\
	-I_{m n} & \mathbf{0}
\end{array}\right)$, where $I_{mn}$ is the identity matrix of order $mn$. 
\end{note}



%%%%%%%%%%%%%%%%%-III-New method of constructing quantum codes--%%%%%%%%%%%%%%%%%%%%%%%%%%%%%%
\section{The necessary and sufficient conditions for QC codes with index even to be symplectic self-orthogonal}\label{sec3}


%Let $g(x)=g_{0}+g_{1} x+g_{2} x+\cdots+ g_{n-1} x^{n-1} \in \mathbb{R}_{q,n}$, $\left[g(x)\right]$ denote vector generated by coefficients of $g(x)$ in $\mathbb{F}_{q}^{n}$, i.e $\left[g(x)\right]=\left[{{g}_{0}},{{g}_{1}},{{g}_{2}},\ldots ,{{g}_{n-1}}\right]$. In adiition, we consider the polynomial $\bar{g}(x)$ defined as $\bar{g}(x)=\sum _{i=0}^{n-1}g_ix^{n-i}$. 
%Moreover, since this paper is studying quasi-cyclic codes with even indices, it is agreed herein that all the indices $\ell$ of quasi-cyclic codes in subsequent texts are even numbers, and let $m=\frac{\ell}{2}$.

In this section, we present the necessary and sufficient conditions for QC codes with index even to be symplectic self-orthogonal. 
First of all, the following lemma is essential for determining the Euclidean inner product between different polynomials in the form of coefficient vectors.


\begin{lemma} \label{commutative_law}(\cite{galindo2018quasi})
	Let $f_a(x)$, $f_b(x)$, and $f_c(x)$ be polynomials in $\mathbb{R}_{q,n}$. Then the following equation holds for the Euclidean inner product among them 
	\begin{equation}
		\left\langle\left[f_a(x) f_b(x)\right],\left[f_c(x)\right]\right\rangle_{e}=\left\langle\left[f_a(x)\right],\left[\overline{f}_b(x) f_c(x)\right]\right\rangle_{e}.
	\end{equation}
\end{lemma}

Since the symplectic inner product can be decomposed into Euclidean inner product, here we determine the Euclidean orthogonality of cyclic codes before discussing the symplectic orthogonality of QC codes. Under the property of cyclic codes, it is easy to derive Theorem \ref{cornerstone}, which establishes necessary and sufficient conditions for Euclidean orthogonality between two cyclic codes.

\begin{theorem}\label{cornerstone}
	Let $g_i(x)$ be a polynomial in $\mathbb{R}_{q,n}$ that divides $x^{n}-1$, and let $\mathscr{C}_i$ be a cyclic code generated by $\left[g_i(x)\right]$, where $i=1$, $2$.
	Then $\mathscr{C}_1$ and $\mathscr{C}_2$ are orthogonal with each other, i.e., for any $a(x)$ and $b(x)$ in $\mathbb{R}_{q,n}$, 
	$\left\langle\left[a(x) g_1(x)\right],\left[b(x) g_2(x)\right]\right\rangle_{e}=0$ holds if and only if $g_1^{\perp_{e}}(x) \mid g_2(x)$.
\end{theorem}


%Next, we will prove the symplectic self-orthogonality of QC codes with index even, beginning with the 1-generator and then the $h$-generator ($h\ge2$). 

\begin{definition}\label{one_quasi-cyclic_def}
	Let $g(x)$ be a polynomial in $\mathbb{R}_{q,n}$ that divides ${{x}^{n}}-1$, and let $f_{j}(x)$ be $\ell$ different polynomials in $\mathbb{R}_{q,n}$ for $0\le j\le \ell-1$, such that the greatest common divisor of $f_0(x),f_1(x),\ldots,f_{\ell-1}(x)$ and $\frac{x^n-1}{g(x)}$ is equal to $1$. 
	%Let $g(x)$ and ${{f}_{j}}(x)$ be polynomials in $\mathbb{R}_{q,n}$, where $g(x)\mid ({{x}^{n}}-1)$, and let $h(x)=\frac{x^n-1}{g(x)}$, $\gcd(f_0(x),f_1(x),\ldots,f_{\ell -1}(x),h(x))=1$. 
	If $\mathscr{C}$ is a QC code generated by $(\left[g(x){f}_{0}(x)\right]$, $\left[g(x){f}_{1}(x)\right],\ldots,\left[g(x){f}_{\ell -1}(x)\right]) \in \mathbb{R}_{q,n}^{\ell} $, then $\mathscr{C}$ is called a $1$-generator QC code with index $\ell$.
	The generator matrix $G$ of $\mathscr{C}$  has the following form:
	\begin{equation}
		G=\left(G_{0}, G_{1}, \ldots, G_{\ell-1}\right) ,
	\end{equation}
	where $G_ {j} $ is an $n\times n$  circulant matrix generated by $\left[g (x) f_{j} (x)\right]$.
\end{definition}

According to \cite{aydin2001structure}, the dimension of $\mathscr{C}$ in Definition \ref{one_quasi-cyclic_def} is given by $n-\deg(g (x))$.
We next present a theorem that determines the necessary and sufficient conditions for $1$-generator QC codes with index even to be symplectic self-orthogonal.

\begin{theorem}\label{one_quasi-cyclic}
	If $\mathscr{C}$ is a $1$-generator QC code in Definition \ref{one_quasi-cyclic_def} with index even $\ell$, then $\mathscr{C}$ is symplectic self-orthogonal if and only if the following equation holds:
	\begin{equation}
		%{{g}^{{{\bot}_{e}}}}(x)\mid \sum\limits_{j=0}^{m-1}{g}(x)({{f}_{j}}(x){{\bar{f}}_{m+j}}(x)-{{f}_{m+j}}(x){{\bar{f}}_{j}}(x)).
		g^{\perp_e}(x)\mid g(x)\Lambda_1,
	\end{equation}
where $\Lambda_1 =\sum\limits_{j = 0}^{m - 1}( {f_j}(x){{\bar f}_{m + j}}(x) - {f_{m + j}}(x){{\bar f}_j}(x) )$ and $m=\frac{\ell}{2}$.
\end{theorem}
\begin{proof}
	Let $a(x)$ and $b(x)$ be arbitrary polynomials in $\mathbb{R}_{q,n}$. Then, any two codewords $\mathbf{c_{1}}$ and $\mathbf{c_{2}}$ of $\mathscr{C}$ can be denoted as $\mathbf{c_{1}}=(\left[a(x) f_{0}(x) g(x)\right],$ $\left[a(x) f_{1}(x) g(x)\right],$ $\ldots, \left[a(x) f_{\ell-1}(x) g(x)\right])$ 
	and $\mathbf{c_{2}}=(\left[b(x) f_{0}(x) g(x)\right],\left[b(x) f_{1}(x) g(x)\right],\ldots, 
	\left[b(x)f_{\ell-1}(x)g(x)\right])$, respectively.
	The symplectic inner product between $\mathbf{c_{1}}$ and $\mathbf{c_{2}}$ is 
		\begin{equation} 
			\begin{array}{rl}
		\left\langle \mathbf{c_{1}}, \mathbf{c_{2}}\right\rangle_{s}&=\mathbf{c_{1}} \cdot \Omega _{mn} \cdot \mathbf{c_{2}}^{T} \\
		&= \sum\limits_{j = 0}^{m - 1} {{{\left\langle {\left[ {a(x)g(x){f_j}(x)} \right],\left[ {b(x)g(x){f_{m + j}}(x)} \right]} \right\rangle }_e}}  \\
		&- \sum\limits_{j = 0}^{m - 1} {{{\left\langle {\left[ {a(x)g(x){f_{m + j}}(x)} \right],\left[ {b(x)g(x){f_j}(x)} \right]} \right\rangle }_e}}  \hfill \\
		&= \sum\limits_{j = 0}^{m - 1} {{{\left\langle {\left[ {a(x)g(x){f_j}(x){{\bar f}_{m + j}}(x)} \right],\left[ {b(x)g(x)} \right]} \right\rangle }_e}}  \\
		&- \sum\limits_{j = 0}^{m - 1} {{{\left\langle {\left[ {a(x)g(x){f_{m + j}}(x){{\bar f}_j}(x)} \right],\left[ {b(x)g(x)} \right]} \right\rangle }_e}}  \hfill \\
		&= \left\langle \left[ a(x) g (x) \Lambda_1\right],\left[b(x)g(x)\right] \right\rangle_e.
	\end{array}	
\end{equation}


	By Theorem \ref{cornerstone}, for all $a(x),b(x)\in\mathbb{R}_{q,n}$, the above equation is equal to zero 
	 if and only if there exists
$g^{\perp e}(x) \mid g (x) \Lambda_1$. 
Thus, we complete the proof.
\end{proof}

%Compared with $1$-generator quasi-cyclic codes, $2$-generator quasi-cyclic codes have a more complex algebraic structure, which can be described as the following definition. 

%\begin{definition} \label{def_2_quasi-cyclic}
%	Let $g_{i}(x)$, $f_{i, j}(x)$ are polynomials in $\mathbb{R}_{q,n}$, and 
%	$g_{i}(x) \mid\left(x^{n}-1\right)$, $1 \leq i \leq 2$, $0 \leq j \leq \ell-1$, then
%	$\left(\left[g_{1}(x) f_{1,0}(x)\right],\left[g_{1}(x) f_{1,1}(x)\right], 
%	\ldots,\left[g_{1}(x) f_{1, \ell-1}(x)\right]\right)$ and
%	$\left(\left[g_{2}(x) f_{2,0}(x)\right],\left[g_{2}(x) f_{2,1}(x)\right], 
%	\ldots, \left[g_{2}(x) f_{2, \ell-1}(x)\right]\right)$ can generate a 2-generator quasi-cyclic code with index $\ell$, whose generator matrix $G$ has the following form:
%	\begin{equation}
%		G=\left(\begin{array}{llll}
%			G_{1,0} & G_{1,1} & \ldots & G_{1, \ell-1} \\
%			G_{2,0} & G_{2,1} & \ldots & G_{2, \ell-1}
%		\end{array}\right),
%	\end{equation}
%	where $G_{i, j}$ are the $n \times n$ circulant matrices generated by $\left[g_{i}(x) f_{i, j}(x)\right]$ respectively.
%\end{definition}
The $h$-generator QC codes can be considered a generalization of $1$-generator QC codes. The definition of $h$-generator QC codes is given here.
\begin{definition}\label{def-h-QC}
	Let  $1 \leq i \leq h$. 
	Let $g_i(x)$ be a polynomial in $\mathbb{R}_{q,n}$ that divides ${{x}^{n}}-1$, and let $f_{i,j}(x)$ be $\ell$ different polynomials in $\mathbb{R}_{q,n}$ for $0\le j\le \ell-1$, such that the greatest common divisor of $f_{i,0}(x),$ $f_{i,1}(x),\ldots,f_{i,{\ell-1}}(x)$ and $\frac{x^n-1}{g_i(x)}$ is equal to $1$.  
	If $\mathscr{C}$ is a QC code with $h$ generators given by:
	$(\left[g_{1}(x) f_{1,0}(x)\right],\left[g_{1}(x) f_{1,1}(x)\right],\ldots,\left[g_{1}(x) f_{1, \ell-1}(x)\right])$, 
	$(\left[g_{2}(x) f_{2,0}(x)\right],$ $\left[g_{2}(x) f_{2,1}(x)\right],\ldots,\left[g_{2}(x) f_{2, \ell-1}(x)\right]),\ldots$, 
	$(\left[g_{h}(x) f_{h, 0}(x)\right],$ $\left[g_{h}(x) f_{h, 1}(x)\right],\ldots,\left[g_{h}(x) f_{h, \ell-1}(x)\right]) \in \mathbb{R}_{q,n}^{\ell} $, then $\mathscr{C}$ is an $h$-generator QC code with index $\ell$, whose generator matrix is
	\begin{equation}\setlength{\arraycolsep}{1.2pt}
		G=\left(\begin{array}{cccc}
			G_{1,0} & G_{1,1} & \ldots & G_{1, \ell-1} \\
			G_{2,0} & G_{2,1} & \ldots & G_{2, \ell-1} \\
			\vdots & \vdots & \ddots & \vdots \\
			G_{h, 0} & G_{h, 1} & \ldots & G_{h, \ell-1}
		\end{array}\right),
	\end{equation}
	where $G_{i, j}$ is an $n \times n$ circulant matrix generated by $\left[g_{i}(x) f_{i, j}(x)\right]$.
\end{definition}

$h$-generator QC codes possess a more complex algebraic structure compared to $1$-generator QC codes. The former can be regarded as a vertical join of $h$ individual $1$-generator QC codes.  The following theorem gives the necessary and sufficient conditions under which the $h$-generator QC codes with index even to be symplectic self-orthogonal.




\begin{theorem}\label{theo-h-QC} 
	If $\mathscr{C}$ is an $h$-generator QC code with index even $\ell$ in Definition \ref{def-h-QC}, then the necessary and sufficient conditions for $\mathscr{C}$ to be symplectic self-orthogonal are that for all $r,s \in\{1,2, \ldots, h\}$, the following equation holds:
 	 \begin{equation}\label{h_sym_quasi-cyclic}
			g_{r}^{\perp_{e}}(x) \mid g_{s}(x)\Lambda_{r,s}, 
	\end{equation} where  $\Lambda_{r,s}=\sum\limits^{m-1}_{i=0} (f_{s, i}(x) \bar{f}_{r, m+i}(x)-f_{s, m+i}(x) \bar{f}_{r, i}(x))$ and $m=\frac{\ell}{2}$.
\end{theorem}
\begin{proof}  
Let $\mathscr{C}_i$ be $1$-generator QC codes generated by $(\left[g_{i}(x) f_{i,0}(x)\right], \left[g_{i}(x) f_{i,1}(x)\right], \ldots, \left[g_{i}(x) f_{i, \ell-1}(x)\right])$, where $i \in \{1,2, \ldots, h\}$.  
Then, in accordance with Definition \ref{def-h-QC}, we have $\mathscr{C} = \mathscr{C}_1 + \mathscr{C}_2 + \cdots + \mathscr{C}_h$.  
  
If $\mathscr{C}$ is a symplectic self-orthogonal code, then for all $\mathbf{c_1}, \mathbf{c_2} \in \mathscr{C}$, $\left\langle \mathbf{c_1}, \mathbf{c_2} \right\rangle_s = 0$. Thus, for all $r, s \in \{1,2, \ldots, h\}$, $\mathscr{C}_r$ and $\mathscr{C}_s$ are symplectic orthogonal to each other.  
  
Conversely, if for all $r, s \in \{1,2, \ldots, h\}$, $\mathscr{C}_r$ and $\mathscr{C}_s$ are symplectic orthogonal to each other, then for all $\mathbf{c_1}, \mathbf{c_2} \in \mathscr{C}$, $\left\langle \mathbf{c_1}, \mathbf{c_2} \right\rangle_s = 0$.   
  
Therefore, $\mathscr{C}$ is a symplectic self-orthogonal code if and only if $\mathscr{C}_r$ and $\mathscr{C}_s$ are symplectic orthogonal to each other for all $r, s \in \{1,2, \ldots, h\}$.  
  
Moreover, using a similar method as the proof of Theorem \ref{one_quasi-cyclic}, we have that $\mathscr{C}_r$ and $\mathscr{C}_s$ are symplectic orthogonal to each other if and only if  
\begin{equation}  
g_{r}^{\perp_{e}}(x) \mid g_{s}(x)\Lambda_{r,s}.  
\end{equation}  
% where  $\Lambda_{r,s}=\sum\limits^{m-1}_{i=0} (f_{s, i}(x) \bar{f}_{r, m+i}(x)-f_{s, m+i}(x) \bar{f}_{r, i}(x))$.
Thus, we complete the proof.  
\end{proof}
%\begin{corollary}\label{theorem_(1,2)-QC}
%	Let the symbols be the same as above, and let $\mathscr{C}$ be a $1$-generator quasi-cyclic code with index $2$ in Definition \ref{one_quasi-cyclic_def}. If it satisfies
%	$g^{\perp_e}(x) \mid g(x)(f_{0}(x) \bar{f}_{1}(x)-f_{1}(x) \bar{f}_{0}(x))$, 
%	then $\mathscr{C}$ is a symplectic self-orthogonal code with parameters $\left[2 n, n-\deg(g(x))\right]_{q}$.
%\end{corollary}
%\begin{proof}
%	By Definition \ref{one_quasi-cyclic_def} and Theorem \ref{one_quasi-cyclic}, this corollary holds.
%	%	In addition, because of $\gcd(f_{j}(x),\frac{x^n-1}{g(x)} )=1$, 
%	%	so $g(x)=\gcd(g(x) f_{0}(x), g(x) f_{1}(x)$, $x^{n}-1)$. 
%	%	Therefore, the dimension of $\mathscr{C}$ is $n-\deg(g(x))$.
%\end{proof}

% \begin{remark}\label{remark1}
% There are three ways to construct symplectic self-orthogonal codes. 
% \begin{itemize}
%     \item The first is to choose appropriate $g_i(x)$ such that $g^{\perp_e}_r(x)\mid g_s(x)$, then search for $f_{i,j}(x)$ to construct symplectic self-orthogonal codes with excellent parameters. 
%     \item The second is to set suitable $f_{i,j}(x)$ such that $f_{s, i}(x) \bar{f}_{r, m+i}(x)-f_{s, m+i}(x) \bar{f}_{r, i}(x) \equiv 0 \pmod{x^n-1}$ holds, which can also produce symplectic self-orthogonal codes. 
%     \item The third is to find suitable $g_i(x)$ and $f_{i,j}(x)$ such that Equation (\ref{h_sym_quasi-cyclic}) holds when $g^{\perp_e}_r(x)\mid g_s(x)$ is not satisfied.
% \end{itemize}
% \end{remark}
As a consequence, for all QC codes with index even, one can directly deduce related necessary and sufficient conditions for symplectic self-orthogonality from Theorems \ref{one_quasi-cyclic} and \ref{theo-h-QC}. 


%\begin{remark}\label{remark2}
%	A special setting of $f(x)$ can be made to satisfy $f(x)=\bar{f}(x)$, which will make the construction process very efficient.
%	Let $u(x) \in \mathbb{R}_{q,\lfloor\frac{n}{2}\rfloor}$, if the structure of $f(x)$ satisfies the following rules:
%	
%	(1) When $n$ is even, let $\left(\left[f(x)\right]\right)=(\left[u(x)\right],\left[\bar{u}(x)\right])$;
%	
%	(2) When $n$ is odd, let $\left(\left[f(x)\right]\right)=(\left[u(x)\right],e,\left[\bar{u}(x)\right])$, where $e\in\mathbb{F}_{q}$;
%	
%	In both cases, $f(x) = \bar{f}(x)$ holds.
%\end{remark}

%
%It is relevant to note that a famous method for constructing QECCs using additive cyclic codes with two generators, 
%is given in the classical work of QECCs \cite{calderbank1998quantum} by Calderbank et al.,  which is shown in the following lemma.
%%given in \cite{calderbank1998quantum}, is similar to our quasi-cyclic construction with index two, as shown in the following lemma.
%\begin{theorem}\label{additive}(\cite{calderbank1998quantum}, Theorem 14)
%	Any quaternary  $(n, 2^{k})$  additive cyclic code  $\mathscr{C}_a$  has two generators, and can be represented as  $\left[\omega p(x)+q(x), r(x)\right]$, where $p(x)$, $q(x)$ and $r(x)$ are binary polynomials, $p(x)$  and  $r(x)$  divide  $x^{n}-1\;(\bmod\; 2)$, $r(x) $ divides  $q(x)(x^{n}- 1)  / p(x)\;(\bmod \;2) $, and  $k=2 n-\deg(p(x))-\deg(r(x)) $. 
%	%If  $\left[\omega p^{\prime}(x)+q^{\prime}(x), r^{\prime}(x)\right]$ is another such representation, then  $p^{\prime}(x)=p(x)$, $r^{\prime}(x)=r(x)$ and  $q^{\prime}(x) \equiv   q(x)(\bmod r(x))$. 
%	Then, $\mathscr{C}$ is trace Hermitian self-orthogonal if and only if
%	\begin{equation}
%	\left\{
%	\begin{array}{ll}
%		p(x) \bar{r}(x) \equiv \bar{p}(x) r(x) \equiv 0 \pmod{ x^{n}-1}, \label{eq-add}\\ 
%p(x) \bar{q}(x) \equiv \bar{p}(x) q(x) \pmod{ x^{n}-1}.  
%	\end{array}
%	\right.
%\end{equation} 
%\end{theorem}
%
%%In fact, there are some quasi-cyclic symplectic self-orthogonal codes in Table \ref{F_QECCs} satisfy Theorem \ref{additive}, such as $\left[\left[105,80,6\right]\right]_2$. However, there are also some QECCs cannot be obtained from Theorem \ref{additive}. We will specify this point below.
%\begin{remark}\label{additive_representation} 
%	Let $\mathscr{C}_a$ denote an additive cyclic code in Theorem \ref{additive}, then $\Phi^{-1}(\mathscr{C}_a)$ is a binary $2$-generator quasi-cyclic code with generators $(p(x),q(x))$ and $(0,r(x))$. In $\mathbb{R}_{2,n}$, polynomials $p(x)$, $q(x)$ and $r(x)$ also can be written as $p(x)=g_1(x)f_0(x)$, $q(x)=g_1(x)f_1(x)$ and $r(x)=g_2(x)f_2(x)$, where $g_1(x)=\gcd(p(x),q(x),x^n-1)$, $g_2(x)=\gcd(r(x),x^n-1)$, $f_0(x)=\frac{p(x)}{g_1(x)}$, $f_1(x)=\frac{q(x)}{g_1(x)}$ and $f_2(x)=\frac{r(x)}{g_2(x)}$. Therefore, quaternary additive cyclic code $\mathscr{C}_a$ also can be represented as $\left[\omega g_1(x)f_0(x)+g_2(x)f_2(x),g_1(x)f_1(x)\right]$.
%\end{remark}

%If $r(x)=0$, then $\Phi^{-1}(\mathscr{C}_a)$ is a $1$-generator quasi-cyclic codes with index two generated by $(g_1(x)f_0(x),g_1(x)f_1(x))$.
%At this time, it is easy to check that Equation (\ref{eq-add}) is equivalent with $g_1^{\perp_{e}}(x) \mid g_1(x)$, which limits the choice of $g_1(x)$. But in Corollary \ref{theorem_(1,2)-QC}, our condition is  $g^{\perp_e}(x) \mid g(x)(f_{0}(x) \bar{f}_{1}(x)-f_{1}(x) \bar{f}_{0}(x))$, one can choose appropriate $f_i(x)$ to make
%$f_{0}(x) \bar{f}_{1}(x)-f_{1}(x) \bar{f}_{0}(x) \equiv 0 \pmod{x^n-1}$ hold, so that all $g(x)$ can be employed to construct symplectic self-orthogonal codes. 
%In this case, Corollary \ref{theorem_(1,2)-QC} cannot be determined by Theorem \ref{additive}.
%Here, we illustrate that with the following example.

\begin{example}  
	Set $q=2$ and $n=40$. Let $g(x)=x^5 + x^4 + x + 1$,   
	$f_0(x)=x^{34} + x^{33} + x^{32} + x^{30} + x^{29} + x^{28} + x^{26} + x^{24} + x^{23} + x^{22} + x^{19} + x^{16} + x^{15} + x^{14} + x^{12} + x^{10} + x^9 + x^8 + x^6 + x^5 + x^4$, and   
	$f_1(x)=x^{37} + x^{35} + x^{34} + x^{30} + x^{29} + x^{23} + x^{20} + x^{18} + x^{15} + x^9 + x^8 + x^4 + x^3 + x$.   
	It is straightforward to verify that $g(x)$, $f_0(x)$ and $f_1(x)$ satisfy the condition in Theorem \ref{one_quasi-cyclic}, namely, $g^{\perp_{e}}(x) \mid g(x)(f_{0}(x) \bar{f}_{1}(x)-f_{1}(x) \bar{f}_{0}(x))$, so we can obtain an $\left[80,35\right]_2^s$ symplectic self-orthogonal code $\mathscr{C}$.   
	Using Magma \cite{bosma1997magma}, we compute that the minimum symplectic dual distance of $\mathscr{C}$ is $10$.  
	  
	According to Theorem \ref{CRSS}, a binary QECC with parameters $\left[\left[40,5,10\right]\right]_2$ can be constructed.    
	It is noted that the best-known binary QECC with length $40$ and dimension $5$ in \cite{Grassltable} has parameters $\left[\left[40,5,9\right]\right]_2$. Thus, the minimum distance record for $\left[\left[40,5\right]\right]_2$ QECCs can be improved to $10$.  
\end{example}


%Therefore, the conditions in Theorem \ref{additive} are not necessary. Here, we supplement it starting from quasi-cyclic codes.

%Then, by Theorem \ref{theo-h-QC}, we can deduce necessary and sufficient conditions for $\mathscr{C}_a$ (resp. $\Phi^{-1}(\mathscr{C}_a)$) to be trace Hermitian (resp. symplectic) self-orthogonal.
%
%\begin{theorem}
%	Let the symbols be the same as Remark \ref{additive_representation}. Then, the quaternary additive cyclic code $\mathscr{C}_a$ (resp. $\Phi^{-1}(\mathscr{C}_a)$) is  trace Hermitian (resp. symplectic) self-orthogonal if and only if 
%	\begin{equation}
%	\left\{
%	\begin{array}{l}
%		g_1^{\perp_e}(x) \mid g_1(x)(f_{0}(x) \bar{f}_{1}(x)-f_{1}(x) \bar{f}_{0}(x)),\\ 
%		g_2^{\perp_e}(x) \mid g_1(x).
%	\end{array}
%	\right.
%\end{equation} 
%\end{theorem}



%%%%%%%%%%%%%%%%%-IV-New method of constructing quantum codes--%%%%%%%%%%%%%%%%%%%%%%%%%%%%%%
\section{Bounds on the minimum symplectic distance of $1$-generator quasi-cyclic codes with index two and their symplectic dual codes}\label{sec4}
%This section describes how to construct $1$-generator and 2-generator quasi-cyclic codes of index two with symplectic self-orthogonality. 
This section presents lower and upper bounds on the minimum symplectic distance of $1$-generator QC codes with index two and their symplectic dual codes. 
In particular, throughout this section, we consider the case where the index $\ell=2$ and $1$-generator QC code $\mathscr{C}$ has generator $(\left[f_0(x)g(x)\right],\left[f_1(x)g(x)\right])$, where $g(x)$ and $f_i(x)$ satisfy the conditions stated in Definition \ref{one_quasi-cyclic_def}. The parity check polynomial of $g(x)$ is denoted as $h(x)=\frac{x^n-1}{g(x)}$. 
%we also construct many symplectic self-orthogonal codes with excellent parameters, for which the associated QECCs are record-breaking. 


%Firstly, in order to decompose quasi-cyclic code, we provide the following definition.
\begin{definition}
	Let $f_a(x)$, $f_b(x)$ be polynomials in $\mathbb{R}_{q,n}$ and $t$ be a positive integer. Define $\mathrm{Circ}_t\left( \left[f_a(x)\right],\left[f_b(x)\right] \right)$ as the following matrix:  
			 %\begin{spacing}{0.1} 
	\begin{equation}
		%\setlength{\arraycolsep}{1.3pt}
		G_t =
		\left(\begin{array}{c|c}
		  \pi(f_a(x))           &  \pi(f_b(x))        \\
			
			\pi(xf_a(x))        & \pi(xf_b(x))        \\
			
			\vdots           & \vdots           \\
			
			\pi(x^{t-1} f_a(x)) & \pi(x^{t-1} f_b(x))
		\end{array}\right)_{t\times 2n}.
	\end{equation}% \end{spacing}
\end{definition}




\begin{lemma}\label{QCmatrix_spread}
	With the above notation.
	Let $\mathscr{C}$ be a $1$-generator QC code with generator $(\left[f_0(x)g(x)\right],\left[f_1(x)g(x)\right])$, and let $t=n-\deg(g(x))-\deg(\gcd( h(x),f_1(x)))-\deg(\gcd( h(x),f_0(x)))$.
	Suppose that $\mathscr{C}_1$, $\mathscr{C}_2$ and $\mathscr{C}_3$ are codes generated by $\left[\frac{x^n-1}{\gcd( h(x),f_1(x))}\right]$, $\left[ \frac{x^n-1}{\gcd( h(x),f_0(x))}\right]$ and $\mathrm{Circ}_t\left( \left[g(x)f_0(x)\right],\left[g(x)f_1(x)\right] \right)$, where $h(x)=\frac{x^n-1}{g(x)}$. Then,  $\mathscr{C}=\mathscr{C}_1\oplus \mathscr{C}_2 + \mathscr{C}_3$.
%	If $g(x)f_0(x)$ and $g(x)f_1(x)$ satisfy $g(x)f_0(x)\mid \frac{x^n-1}{\gcd(h(x),f_1(x))} \pmod{x^n-1}$ and $g(x)f_0(x)\mid \frac{x^n-1}{\gcd(h(x),f_1(x))} \pmod{x^n-1}$ .
	%there exist polynomials $r_1(x)$,$r_2(x)$ and $r_3(x)$ in $\mathbb{R}_{q,n}$ such that a generator matrix $G$ of the 
\end{lemma}
\begin{proof}
Since the generator of $\mathscr{C}$ is $(\left[f_0(x)g(x)\right],\left[f_1(x)g(x)\right])$, any codeword of $\mathscr{C}$ has the form  $\mathbf{c}=(\mathbf{c}^{\prime},\mathbf{c}^{{\prime}{\prime}})=(\left[a(x)g(x)f_0(x)\right],$ $\left[a(x)g(x)f_1(x)\right])$, where $a(x)\in \mathbb{R}_{q,n}$. 
	This implies that only the following three types of nonzero codewords may appear in $\mathscr{C}$. We represent them with three sets:
	\begin{equation}
		\begin{matrix}
			S_1=\{ (\mathbf{c}^{\prime},\mathbf{c}^{{\prime}{\prime}})\mid \mathbf{c}^{\prime}\ne\mathbf{0}_n , \mathbf{c}^{{\prime}{\prime}} =\mathbf{0}_n \}, \\
			S_2=\{ (\mathbf{c}^{\prime},\mathbf{c}^{{\prime}{\prime}})\mid \mathbf{c}^{\prime}=\mathbf{0}_n , \mathbf{c}^{{\prime}{\prime}}\ne \mathbf{0}_n\}, \\
			S_3=\{ (\mathbf{c}^{\prime},\mathbf{c}^{{\prime}{\prime}})\mid \mathbf{c}^{\prime}\ne\mathbf{0}_n , \mathbf{c}^{{\prime}{\prime}}\ne \mathbf{0}_n\}.
		\end{matrix}
	\end{equation}


	If $\mathbf{c} \in S_1$, then $a(x)g(x)f_0(x) \not\equiv 0 \pmod{x^n-1}$ and $a(x)g(x)f_1(x)\equiv 0 \pmod{x^n-1}$. We have $\frac{x^n-1}{\gcd(x^n-1,g(x)f_1(x))}\mid a(x)$ and $a(x)$ can be written as $a(x)\equiv r(x)\frac{x^n-1}{\gcd(x^n-1,g(x)f_1(x))} \pmod{x^n-1}$, where $r(x)$ is a polynomial in $\mathbb{R}_{q,n}$ of degree less than $\deg(\gcd(x^n-1,g(x)f_1(x)))$. Since $\gcd(f_0(x),f_1(x),h(x))=1$ and $\frac{x^n-1}{\gcd(x^n-1,g(x)f_1(x))}g(x)=\frac{x^n-1}{\gcd(h(x),f_1(x))}$, we get $a(x)g(x)f_0(x) \equiv 0 \pmod{\frac{x^n-1}{\gcd(h(x),f_1(x))}}$.
 At this time, the cyclic code determined by $\left[a(x)g(x)f_0(x)\right]$ is indeed $\mathscr{C}_1$, which has dimension $\deg(\gcd( h(x),f_1(x)))$.
	
	Similarly, if $\mathbf{c} \in S_2$, then $a(x)g(x)f_0(x)\equiv 0 \pmod{x^n-1}$ and $a(x)g(x)f_1(x) \not\equiv 0 \pmod{x^n-1}$. 
 Therefore, when running through all $a(x)$ that divide $\frac{x^n-1}{\gcd(x^n-1,g(x)f_0(x))}$, the cyclic code that composed of all codewords $\left[a(x)g(x)f_0(x)\right]$ is $\mathscr{C}_2$, whose dimension is $\deg(\gcd( h(x),f_0(x)))$.
	Therefore, $\mathscr{C}_1\oplus \mathscr{C}_2 \subset \mathscr{C}$ and $\dim(\mathscr{C}_1\oplus \mathscr{C}_2)=\deg(\gcd( h(x),f_1(x)))+\deg(\gcd( h(x),f_0(x)))$.
	
	Given that $t=n-\deg(g(x))-\deg(\gcd( h(x),f_1(x)))-\deg(\gcd( h(x),f_0(x)))$, the nonzero codewords generated by $\mathrm{Circ}_t\left( \left[g(x)f_0(x)\right],\left[g(x)f_1(x)\right] \right)$ are all contained in $S_3$, and the dimension of $\mathscr{C}_3$ is $t$. Therefore, we have $\mathscr{C}_3\subset \mathscr{C}$ and $\dim(\mathscr{C}_1\oplus \mathscr{C}_2+\mathscr{C}_3)=n-\deg(g(x))$. 
	Therefore, the generator matrix of $\mathscr{C}$ can be written in the following form:
		\begin{equation}\label{QCtransform}
	%\setlength{\arraycolsep}{1.2pt}
	G^{\prime}=
	%		\left(\begin{array}{c}
		%			G_1\\\hline
		%			G_2\\\hline
		%			G_3
		%		\end{array}\right)=
	\left(\begin{array}{c|c}
		\pi( g(x)f_0(x))                                                                  & \pi( g(x)f_1(x))                                                             \\
		
		\pi(x g(x)f_0(x))                                                                 & \pi(x g(x)f_1(x))                                                            \\
		
		\vdots                                                                           & \vdots                                                                      \\
		
		\pi(x^{t-1} g(x)f_0(x))         & \pi(x^{t-1}g(x)f_1(x))     \\
		\hline
		\mathbf{0}_n                                                                     & \pi( \frac{x^n-1}{\gcd( h(x),f_0(x))})                                  \\
		\mathbf{0}_n                                                                     & \pi(x \frac{x^n-1}{\gcd( h(x),f_0(x))} )                                \\
		\vdots                                                                           & \vdots                                                                      \\
		\mathbf{0}_n                                                                     & \pi(x^{t_a-1}   \frac{x^n-1}{\gcd( h(x),f_0(x))})\\		\hline
  		\pi( \frac{x^n-1}{\gcd( h(x),f_1(x))})                                 & \mathbf{0}_n                                                                \\
		
		\pi(x  \frac{x^n-1}{\gcd( h(x),f_1(x))})                               & \mathbf{0}_n                                                                \\
		
		\vdots                                                                           & \vdots                                                                      \\
		
		\pi(x^{t_b-1}  \frac{x^n-1}{\gcd( h(x),f_1(x))}) & \mathbf{0}_n                                                                \\

	\end{array}\right),\\%前两块矩阵
\end{equation}
	where $t_a=\deg(\gcd( h(x),f_1(x)))$ and $t_b=\deg(\gcd( h(x),f_0(x)))$. 
\end{proof}

%//-----------------------------------------------------------------------------\textbf{delete?}
%\begin{lemma}\label{HQCBound}
%	Let $\mathscr{C}$ be a $1$-generator quasi-cyclic code in Definition \ref{one_quasi-cyclic_def},  and $\ell=2$.
%	Suppose that $C_1$, $C_2$, $C_3$ and $C_4$ are cyclic codes determined by $\left[ g(x)f_0(x)\right]$, $\left[ g(x)f_1(x)\right]$, $\left[ \frac{x^n-1}{\gcd(h(x),f_1(x))} \right]$ and $\left[ \frac{x^n-1}{\gcd(h(x),f_0(x))}\right] $, respectively. Then, the quasi-cyclic code $\mathscr{C}$ determined by $(\left[g(x)f_0(x)\right],\left[g(x)f_1(x)\right])$ satisfy 
%		\begin{equation}
%		\min\left\{ d_H(C_3),d_H(C_4), d_H(C_1)+d_H(C_2) \right\}  \le d_H(\mathscr{C}) \le 
%		\min\left\{ d_H(C_3),d_H(C_4)\right\}
%	\end{equation}
%%	\begin{equation} 
%%		\begin{aligned}
%%				\min\left\{  { d_H\left(\frac{x^n-1}{\gcd(h(x),f_1(x))}\right) , d_H\left(\frac{x^n-1}{\gcd(h(x),f_0(x))}\right),} 
%%				\right. \\
%%				\left.
%%				 {\frac{d_H(g(x)f_0(x))+d_H(g(x)f_1(x))+(q-1)d_H(\gcd(g(x)f_0(x),g(x)f_1(x)))}{q} } \right\}  \le d_s(\mathscr{C}) 
%%		\end{aligned}
%%\end{equation}
%\end{lemma}
%\begin{proof}
%From Lemma \ref{QCmatrix_spread}, it is easily deduce that this Lemma holds.
%\end{proof}
%
%\begin{example}
%	Let $q=2$ and $n=15$. Set $g(x)=1$, $f_0(x)=x^5 + x^4 + x^2 + 1$ and $f_1(x)=x^4 + x^3 + 1$. Since $\gcd(f_0(x),f_1(x))=1$, $(\left[g(x)f_0(x)\right],\left[g(x)f_1(x)\right])$ will generate a $\left[30,15\right]_2$ quasi-cyclic code $\mathscr{C}$.
%	One can easily check that the distances of cyclic codes generated by $g(x)f_0(x)$, $g(x)f_1(x)$,  $f_0(x)\frac{x^n-1}{gcd( x^n-1,f_1(x))}$ and $\frac{x^n-1}{\gcd(h(x),f_0(x))}$ are $4$, $3$, $8$ and $7$, respectively. By Lemma \ref{HQCBound}, 
%	we can immediately conclude that the lower bound on the minimum distance of $\mathscr{C}$ is $7$, as is the upper bound. 
%	Therefore, the resulting quasi-cyclic code $\mathscr{C}$ is a $\left[30, 15, 7\right]_2$ code.
%	Further, using Magma \cite{bosma1997magma}, we can compute that weight distribution of $\mathscr{C}$ is $1+60z^7+165z^8+560z^9+1200z^{10}+\cdots+15z^{24}$.
%\end{example}
%//-----------------------------------------------------------------------------




 \begin{theorem}\label{SQCBound} 
	With the above notation,
	let $\mathscr{C}$ be a $1$-generator QC code with generator $(\left[f_0(x)g(x)\right],\left[f_1(x)g(x)\right])$.
 For $\alpha\in \mathbb{F}_q^*$, define
$\mathcal{I}^{(\alpha)}(\mathscr{C})=\gcd(f_0(x)+\alpha f_1(x),h(x))$. 
% $\mathcal{I}^{(\alpha)}(\mathscr{C})=\gcd(f_0(x)+\alpha f_1(x),h(x))$. 
 Let $\mathcal{S} =\{\mathcal{I}^{(\alpha)}(\mathscr{C}) \mid \mathcal{I}^{(\alpha)}(\mathscr{C})\ne 1, \forall\alpha\in\mathbb{F}_q^*\}$.
	Suppose that $\mathscr{C}_0$, $\mathscr{C}_1$, $\mathscr{C}_2$, $\mathscr{C}_3$, $\mathscr{C}_4$ and $\mathscr{C}_5$ are cyclic codes determined by $[g(x)]$, $[g(x)f_0(x)]$, $[g(x)f_1(x)]$, $ \left[ \frac{x^n-1}{\gcd(h(x),f_0(x))}\right]$, $\left[\frac{x^n-1}{\gcd(h(x),f_1(x))}\right]$ and $\left[ g(x)\operatorname{lcm}(f_0(x),f_1(x))\right]$, respectively.
 Define $\mathcal{D}(\mathscr{C})$ as in Equation (\ref{The. 4. E.1}). \begin{figure*}
\begin{equation}\label{The. 4. E.1}
\mathcal{D}(\mathscr{C})=\max \left\{ \left\lceil\left(d_H(\mathscr{C}_1)+d_H(\mathscr{C}_2)+(q-|\mathcal{S} |-1)d_H(\mathscr{C}_0)+\sum\limits_{\mathcal{I}^{(\alpha)}(\mathscr{C})\in \mathcal{S} } d_H\left(\left[ g(x)\mathcal{I}^{(\alpha)}(\mathscr{C}) \right]\right)\right)/ q\right\rceil, d_H(\mathscr{C}_1),d_H(\mathscr{C}_2)\right\}.
\end{equation} 
\end{figure*}
Writing $d_{\mathrm{upper}}=	\min\left\{d_H(\mathscr{C}_3),d_H(\mathscr{C}_4)\right\}$ and   $d_{\mathrm{lower}}$ as in Equation (\ref{symplectic_bound}),
\begin{figure*}\begin{equation}\label{symplectic_bound}
	d_{\mathrm{lower}}=
	\left\{
	\begin{array}{ll}
		\min \left\{ d_H(\mathscr{C}_3),d_H(\mathscr{C}_4),\mathcal{D}(\mathscr{C}) \right\},& \mathcal{S} =\emptyset,                                                           \\
		\min\left\{ d_H(\mathscr{C}_3),d_H(\mathscr{C}_4), d_H(\mathscr{C}_5), \mathcal{D}(\mathscr{C})\right\},&q=2,\mathcal{S}\ne \emptyset,                     \\
	\min \left\{ d_H(\mathscr{C}_3),d_H(\mathscr{C}_4),\max\{d_H(\mathscr{C}_1),d_H(\mathscr{C}_2)\} \right\},&q>2, \mathcal{S}\ne \emptyset.                                        \\
	\end{array}
	\right.
\end{equation} \end{figure*}
then we have $d_{\mathrm{lower}}\le d_s(\mathscr{C})\le d_{\mathrm{upper}}$.
\end{theorem}
\begin{proof}
 Denote codeword of $\mathscr{C}$ by $\mathbf{c}=(\mathbf{c}^{\prime},\mathbf{c}^{\prime\prime})=(\left[a(x)g(x)f_0(x)\right],\left[a(x)g(x)f_1(x)\right])$, where $a(x)\in \mathbb{R}_{q,n}$. 
According to \cite{ling2010generalization}, we have
$q\mathbf{w}_{s}(\mathbf{c})=\mathbf{w}_{H}(\mathbf{c}^{\prime})+\mathbf{w}_{H}(\mathbf{c}^{\prime\prime})+\sum\limits_{\alpha\in \mathbb{F}_q^*} \mathbf{w}_{H}(\mathbf{c}^{\prime}+\alpha \mathbf{c}^{\prime\prime})$.


If $\mathbf{c}^{\prime}=\mathbf{0}_n$ or $\mathbf{c}^{\prime\prime}=\mathbf{0}_n$, then we have $a(x)g(x)f_0(x)\equiv  0 \pmod{x^n-1}$ or $a(x)g(x)f_1(x) \equiv  0 \pmod{x^n-1}$.
By Lemma \ref{QCmatrix_spread}, the nonzero half part of $\mathbf{c}$ is a codeword of $\mathscr{C}_3$ or $\mathscr{C}_4$. 
Therefore, it follows that 
 $d_s(\mathscr{C})\le \min \{d_H(\mathscr{C}_3), d_H(\mathscr{C}_4)\}$. 


If $\mathbf{c}^{\prime}\ne \mathbf{0}_n$ and $\mathbf{c}^{\prime\prime}\ne \mathbf{0}_n$, then $a(x)g(x)f_0(x)\not\equiv 0 \pmod{x^n-1}$ and $a(x)g(x)f_1(x) \not\equiv 0 \pmod{x^n-1}$.
There are two cases to consider. 
Firstly, if for all $ \alpha$ in $\mathbb{F}_q^*$, $\gcd(f_0(x)+\alpha f_1(x),h(x))= 1$, i.e., $\mathcal{S} =\emptyset$,
then the following formula holds.
\begin{equation}\label{ESBound1}
	\begin{array}{rl}
		d_s(\mathscr{C})\ge&  (d_H(\left[ g(x)f_0(x)\right])+d_H(\left[g(x)f_1(x)\right])\\
		&+(q-1) d_H(\left[ g(x)\right] ))/q.
	\end{array}
\end{equation}


Secondly, if there is $\alpha$ in $\mathbb{F}_q^*$ such that $\mathcal{I}^{(\alpha)}(\mathscr{C})=\gcd(f_0(x)+\alpha f_1(x),h(x))\ne 1$, i.e., $\mathcal{S} \ne \emptyset$, then there exist  $a(x)$ of degree less than $n-\deg(g(x))$ such that $a(x)(f_0(x)+\alpha f_1(x)) g(x)\equiv 0 \pmod{x^n-1}$, i.e.,  $\mathrm{Supp}(\mathbf{c}^{\prime})=\mathrm{Supp}(\mathbf{c}^{\prime\prime})$. 
For $q=2$, if $\mathrm{Supp}(\mathbf{c}^{\prime})=\mathrm{Supp}(\mathbf{c}^{\prime\prime})$, then we have $\mathbf{c}^{\prime}=\mathbf{c}^{\prime\prime}$ $\in \mathscr{C}_1 \cap \mathscr{C}_2$. Therefore, 
$d_s(\mathscr{C}) \ge d_H( \left[ g(x) \operatorname{lcm}(f_0(x),f_1(x)) \right] )$.  
For $q> 2$, we can only determine $d_s(\mathscr{C})\ge \max\{d_H(\mathscr{C}_1),d_H(\mathscr{C}_2)\}$.


If $\mathrm{Supp}(\mathbf{c}^{\prime})\ne \mathrm{Supp}(\mathbf{c}^{\prime\prime})$, then we have Equation (\ref{ESBound2}). 
\begin{figure*}\begin{equation}\label{ESBound2}
d_s(\mathscr{C})\ge \left(d_H(\mathscr{C}_1)+d_H(\mathscr{C}_2)+(q-|\mathcal{S} |-1)d_H(\mathscr{C}_0)+\sum\limits_{\mathcal{I}^{(\alpha)}(\mathscr{C})\in \mathcal{S} } d_H\left(\left[ g(x)\mathcal{I}^{(\alpha)}(\mathscr{C}) \right]\right)\right)/ q.
\end{equation}\end{figure*}

Additionally, $d_s(\mathscr{C})$ must also satisfy the condition that it is greater than or equal to $\max\left\{ d_H(\mathscr{C}_1),d_H(\mathscr{C}_2) \right\}$ when Equations (\ref{ESBound1}) and (\ref{ESBound2}) are satisfied. 
Therefore, we have $d_s(\mathscr{C})\ge\mathcal{D}(\mathscr{C})$.

From the combination of the above conditions one can obtain upper and lower bounds for the minimum symplectic distance of $\mathscr{C}$. This concludes the proof.
\end{proof}

%\begin{example}\label{good_Ex1}
%	Let $q=2$ and $n=21$. Set $g(x)=1$, $f_0(x)=x^{20} + x^{18} + x^{13} + x^{12} + x^{11} + x^{10} + x^9 + x^8 + x^7 + x^2$ and $f_1(x)=x^{20} + x^{19} + x^{18} + x^{17} + x^{13} + x^{12} + x^9 + x^6 + x^3 + x^2 + 1$. Since $\deg(g(x))=0$, $(\left[g(x)f_0(x)\right],\left[g(x)f_1(x)\right])$ will generate a $\left[42,21\right]_2$ quasi-cyclic code $\mathscr{C}$.
%	One can easily check that $f_{\alpha}(x)=\gcd(f_0(x)+f_1(x),x^{21}-1)=x^2 + x + 1$ and the distances of cyclic codes generated by $g(x)f_0(x)$, $g(x)f_1(x)$,  $f_0(x)\frac{x^n-1}{gcd( x^n-1,f_1(x))}$, $\frac{x^n-1}{\gcd(h(x),f_0(x))}$, $g(x)\operatorname{lcm}(f_0(x),f_1(x))$ and $g(x)f_{\alpha}(x)$ are $2$, $5$, $21$, $8$, $6$ and $2$, respectively. By Theorem \ref{SQCBound}, 
%    it follows that the lower bound on the minimum symplectic distance of $\mathscr{C}$ is $5$, while the upper bound is $8$.
%	Using Magma \cite{bosma1997magma}, we can compute that the resulting quasi-cyclic code $\mathscr{C}$ has a real symplectic distance $8$. Thus, $\mathscr{C}$ corresponds to an additive cyclic $(21, 10.5, 8)_4$ code, which outperforms the best quaternary linear $\left[21, 10, 8\right]_4$ code.
%\end{example}

\begin{remark}\label{improve_bound}
	Since $1$-generator QC codes with index two over $\mathbb{F}_q$ are equivalent with additive cyclic codes over $\mathbb{F}_{q^2}$ \cite{GuanAdditive2023}, Theorem \ref{SQCBound} also can be viewed as bounds on the minimum Hamming distance on additive cyclic codes.
	In addition, Corollary 3.9 in \cite{dastbasteh2023polynomial} can be viewed as a subcase ($q>2,\mathcal{S} \ne\emptyset$) of Theorem \ref{SQCBound},  
	and Theorem 1 in \cite{GuanAdditive2023} is a simplified form of $\mathcal{S} =\emptyset$ in Theorem \ref{SQCBound}.
	Thus, Theorem \ref{SQCBound} improves and enhances the results in \cite{dastbasteh2023polynomial} and \cite{GuanAdditive2023}. 
	%, in particular,  Corollary 3.9 in \cite{dastbasteh2023polynomial} can be viewed as $\mathcal{S} \ne\emptyset$ in Equation (\ref{symplectic_bound}), and Theorem 1 in \cite{GuanAdditive2023} is a subcase of $\mathcal{S} =\emptyset$ in Equation (\ref{symplectic_bound}).
	%	In fact, Corollary 3.9 in \cite{dastbasteh2023polynomial} and Theorem 1 in \cite{GuanAdditive2023} are subcases of Theorem \ref{SQCBound}. Corollary 3.9 in \cite{dastbasteh2023polynomial} can be viewed as $\mathcal{S} \ne\emptyset$ in Equation (\ref{symplectic_bound}), and Theorem 1 in \cite{GuanAdditive2023} is $\mathcal{S} =\emptyset$ in Equation (\ref{symplectic_bound}).
\end{remark}


\begin{example}\label{good_Ex1}
	Take $q=2$ and $n=21$. Set $g(x)=x^6 + x^5 + x^4 + x^2 + 1$, $f_0(x)=x^{11} + x^6 + x^5 + x^2 + x + 1$ and $f_1(x)=x^{12} + x^7 + x^6 + x^5 + x^2 + x + 1$, then $h(x)=\frac{x^{21}-1}{g(x)}=x^{15} + x^{14} + x^{12} + x^9 + x^8 + x^5 + x^2 + 1$. Since $\gcd(f_0(x),f_1(x),h(x))=1$ and $\deg(h(x))=15$, $([g(x)f_0(x)],[g(x)f_1(x)])$ generates a $\left[42,15\right]_2$ QC code $\mathscr{C}$.


 Given that $\mathcal{S}=\{ \mathcal{I}^{(1)}(\mathscr{C})=\gcd(f_0(x)+f_1(x),h(x))=x^5 + x^4 + 1\}$, the lower bound on the minimum symplectic distance of $\mathscr{C}$ is the second condition in Equation (\ref{symplectic_bound}), i.e., 
 \begin{equation}
	\begin{array}{rl}
	 \min\left\{ d_H(\mathscr{C}_3),d_H(\mathscr{C}_4), d_H(\mathscr{C}_5), \mathcal{D}(\mathscr{C}) \right\}  \le d_s(\mathscr{C}) &\\
	 \le\min\left\{d_H(\mathscr{C}_3),d_H(\mathscr{C}_4)\right\}&
	\end{array}
\end{equation}
% $$ \min\left\{ d_H(\mathscr{C}_3),d_H(\mathscr{C}_4), d_H(\mathscr{C}_5), \mathcal{D}(\mathscr{C}) \right\} \le d_s(\mathscr{C})\le \min\left\{d_H(\mathscr{C}_3),d_H(\mathscr{C}_4)\right\},$$
  where $\mathcal{D}(\mathscr{C})=\max \left\{ \left\lceil\left(d_H(\mathscr{C}_1)+d_H(\mathscr{C}_2)+d_H(\left[g(x)\mathcal{I}^{(1)}(\mathscr{C})\right])\right)/ 2\right\rceil\right.$, 
  $\left.d_H(\mathscr{C}_1),d_H(\mathscr{C}_2)\right\}$. 
Here we list the generator polynomials and parameters of the cyclic codes $\mathscr{C}_i$:
 \begin{itemize}
      \item  $\mathscr{C}_0$: Generator polynomial $\left[g(x)\right]$. Parameters $\left[21,15,3\right]_2$.
      \item  $\mathscr{C}_1$: Generator polynomial $\left[g(x)f_0(x)\right]$. Parameters $\left[21,14,4\right]_2$.
      \item  $\mathscr{C}_2$: Generator polynomial $\left[g(x)f_1(x)\right]$. Parameters $\left[21,6,7\right]_2$.
     \item $\mathscr{C}_3$: Generator polynomial $\left[\frac{x^{21}-1}{\gcd(h(x),f_0(x))}\right]=\left[\frac{x^{21}-1}{x+1}\right]$. Parameters $\left[21,1,21\right]_2$.
      \item $\mathscr{C}_4$: Generator polynomial $\left[\frac{x^{21}-1}{\gcd(h(x),f_1(x))}\right]=\left[x^{12} + x^{11} + x^9 + x^7 + x^3 + x^2 + x + 1\right]$. Parameters $\left[21,9,8\right]_2$.
     \item $\mathscr{C}_5$: Generator polynomial $\left[g(x)\operatorname{lcm}(f_0(x),f_1(x))\right]=\left[x^{16} + x^{15} + x^{14} + x^{13} + x^{12} + x^{10} + x^8 + x^5 + x^4 + 1\right]$. Parameters $\left[21,5,10\right]_2$.
     \item   Generator polynomial $\left[g(x)\mathcal{I}^{(1)}(\mathscr{C})\right]=\left[x^{11} + x^8 + x^7 + x^2 + 1\right]$. Parameters $\left[21,10,5\right]_2$. 
 \end{itemize}
 
 
% 	 $\min\left\{ d_H\left(\left[\right]\right),d_H\left(\left[\frac{x^{21}-1}{\gcd(h(x),f_1(x))}\right]\right), d_H\left(\left[ g(x)\operatorname{lcm}(f_0(x),f_1(x))\right]\right), \mathcal{D}_{\perp}(\mathscr{C})\right\}$, where
  
% 	 $\mathcal{D}_{\perp}(\mathscr{C})=\max \left\{ \left\lceil\left(d_H(\left[ g(x)f_0(x) \right])+d_H(\left[ g(x)f_1(x)\right])+d_H(\left[ g(x)\mathcal{I}^{(1)}(\mathscr{C}) \right])\right)/ 2\right\rceil, d_H(\left[ g(x)f_0(x) \right]),d_H(\left[ g(x)f_1(x)\right])\right\}$,
% and the upper bound is $\min\left\{d_H\left(\left[\frac{x^{21}-1}{\gcd(h(x),f_0(x))}\right]\right),d_H\left(\left[\frac{x^{21}-1}{\gcd(h(x),f_1(x))}\right]\right)\right\}$.

We can directly derive that the lower and upper bounds on the minimum symplectic distance of $\mathscr{C}$ are both $8$ based on the parameters of the above cyclic codes. 
 By utilising Magma \cite{bosma1997magma}, it can be determined that the specific parameters of $\mathscr{C}$ are indeed $\left[42,15,8\right]_2^s$. 
 In contrast, Corollary 3.9 in \cite{dastbasteh2023polynomial} provides a lower bound of $4$ on the minimum symplectic distance of $\mathscr{C}$, whereas Theorem 1 in \cite{GuanAdditive2023} does not yield a lower bound for this code. Consequently, our bound is more sharper and applicable.
\end{example}

The following example shows that Theorem \ref{SQCBound} performs better even though the constraints in Theorem 1 of \cite{GuanAdditive2023} are satisfied.

\begin{example}\label{good_Ex_for additive}
	Let $q=2$ and $n=31$. Set $g(x)=x^5 + x^2 + 1$, $f_0(x)=x^{25} + x^{24} + x^{23} + x^{22} + x^{21} + x^{20} + x^{19} + x^{16} + x^{13} + x^{11} + x^9 +x^8 + x^7 + x^6 + 1$ and $f_1(x)=x^{21} + x^{20} + x^{15} + x^{13} + x^8 + x^5 + x^4 + x^3$, then $h(x)=\frac{x^{31}-1}{g(x)}=x^{26} + x^{23} + x^{21} + x^{20} + x^{17} + x^{16} + x^{15} + x^{14} + x^{13} + x^9 + x^8 + x^6 +
	x^5 + x^4 + x^2 + 1$. Since $\gcd(f_0(x),f_1(x),h(x))=1$ and $\deg(h(x))=26$, $([g(x)f_0(x)],[g(x)f_1(x)])$ produces a $\left[62,26\right]_2$ QC code $\mathscr{C}$.
 

 Given that $\mathcal{S}=\emptyset$, the lower bound on the minimum symplectic distance of $\mathscr{C}$ is the first condition in Equation (\ref{symplectic_bound}), i.e., 
  \begin{equation}
 	\begin{array}{rl}
 		\min\left\{ d_H(\mathscr{C}_3),d_H(\mathscr{C}_4), \mathcal{D}(\mathscr{C})\right\} \le d_s(\mathscr{C}) &\\
 		\le \min\left\{d_H(\mathscr{C}_3),d_H(\mathscr{C}_4)\right\}&
 	\end{array}
 \end{equation}
% $$\min\left\{ d_H(\mathscr{C}_3),d_H(\mathscr{C}_4), \mathcal{D}(\mathscr{C})\right\} \le d_s(\mathscr{C})\le \min\left\{d_H(\mathscr{C}_3),d_H(\mathscr{C}_4)\right\},$$  
 where $\mathcal{D}(\mathscr{C})=\max \left\{ \left\lceil \left(d_H(\mathscr{C}_0)+d_H(\mathscr{C}_1)+d_H(\mathscr{C}_2)\right)/2\right\rceil \right.,$ $\left.d_H(\mathscr{C}_1),d_H(\mathscr{C}_2)\right\}$.
Here we list the generator polynomials and parameters of the cyclic codes $\mathscr{C}_i$:
 \begin{itemize}
      \item  $\mathscr{C}_0$: Generator polynomial $\left[g(x)\right]$. Parameters $\left[31,26,3\right]_2$.
      \item  $\mathscr{C}_1$: Generator polynomial $\left[g(x)f_0(x)\right]$. Parameters $\left[31,16,7\right]_2$.
      \item  $\mathscr{C}_2$: Generator polynomial $\left[g(x)f_1(x)\right]$. Parameters $\left[31,25,4\right]_2$.
     \item $\mathscr{C}_3$: Generator polynomial $\left[\frac{x^{31}-1}{\gcd(h(x),f_0(x))}\right]=\left[x^{21} +\right. $ $ \left. x^{18} + x^{17} + x^{15} +x^{13} + x^{10} + x^5 + x^4 + x^3 + x^2 + x \right. $ $ \left.+ 1\right]$. Parameters $\left[31,10,12\right]_2$.
      \item $\mathscr{C}_4$: Generator polynomial $\left[\frac{x^{31}-1}{\gcd(h(x),f_1(x))}\right]=\left[\frac{x^{31}-1}{x+1}\right]$. Parameters $\left[31,1,31\right]_2$.
 \end{itemize}

	% It is easy to check that $\mathcal{S}=\emptyset$, so by Theorem \ref{SQCBound}, the lower bound on the minimum symplectic distance of $\mathscr{C}$ is
	% $\min\left\{ d_H\left(\left[\frac{x^{31}-1}{\gcd(h(x),f_0(x))}\right]\right),d_H\left(\left[\frac{x^{31}-1}{\gcd(h(x),f_1(x))}\right]\right), \mathcal{D}_{\perp}(\mathscr{C})\right\}$, 
	% where
	% $\mathcal{D}_{\perp}(\mathscr{C})=\max \left\{ \left\lceil\left(d_H(\left[ g(x)f_0(x) \right])+d_H(\left[ g(x)f_1(x)\right])+d_H(\left[ g(x) \right])\right)/ 2\right\rceil, d_H(\left[ g(x)f_0(x) \right]),d_H(\left[ g(x)f_1(x)\right])\right\}$,
	% and the upper bound is $\min\left\{d_H\left(\left[\frac{x^{31}-1}{\gcd(h(x),f_0(x))}\right]\right),d_H\left(\left[\frac{x^{31}-1}{\gcd(h(x),f_1(x))}\right]\right)\right\}$.
	
	We can directly derive lower and upper bounds on the minimum symplectic distance of $\mathscr{C}$ are $7$ and $12$ based on the parameters of above cyclic codes. By using Magma \cite{bosma1997magma}, we can calculate that the specific parameters of $\mathscr{C}$ are $\left[62,26,11\right]_2^s$. 
 In contrast, Theorem 1 in \cite{GuanAdditive2023} provides a lower bound of $5$ on the minimum symplectic distance of $\mathscr{C}$. Hence, our bound is more tighter. 
\end{example}



\begin{theorem}\label{Dual_Structure} 
	With the above notation,
	let $\mathscr{C}$ be a $1$-generator QC code with generator $(\left[f_0(x)g(x)\right],\left[f_1(x)g(x)\right])$.
	 If $\gcd(f_0(x),g(x))=1$,		
then the symplectic dual code $\mathscr{C}^{\perp_s}$ of $\mathscr{C}$ is a 2-generator QC code, whose generator matrix is
\begin{equation}
	G^{\perp_{s}}=
	\begin{pmatrix}
		F_{0}	& F_{1}\\
		\mathbf{0}&G_{g^{\perp_e}}
	\end{pmatrix},
\end{equation}
where $(F_{0},F_1)$ and $(\mathbf{0},G_{g^{\perp_e}})$ are the generator matrices of $1$-generator QC codes determined by $(\left[ \bar{f}_{0}(x)\right], \left[ \bar{f}_{1}(x)\right])$ and $(\left[ 0\right], \left[ g^{\perp_e}(x)\right])$, respectively. 
 %where $(F_{0},F_1)$ is an $n \times 2n$ circulant matrices generated by $(\left[ \bar{f}_{0}(x)\right], \left[\bar{f}_{1}(x)\right])$, $G_{g^{\perp_e}}$ is the generator matrix of the cyclic code generated by $\left[g^{\perp_e}(x)\right]$, and $\mathbf{0}$ is an appropriate zero matrix.
\end{theorem}
\begin{proof}
	 Let $\mathscr{C}^{\circ}$ denote the code generated by $G^{\perp_{s}}$.  
 Then any codeword of $\mathscr{C}$ and $\mathscr{C}^{\circ}$ has the form $\mathbf{c_{1}}=(\left[a(x) f_{0}(x) g(x)\right],$ $\left[a(x) f_{1}(x) g(x)\right])$ and $\mathbf{c_{2}}=(\left[b(x) \bar{f}_{0}(x)\right],\left[b(x) \bar{f}_{1}(x)+c(x)g^{\perp_e}(x)\right])$, where $a(x),b(x),c(x)\in\mathbb{R}_{q,n}$. 
	Using a similar approach of Theorem \ref{one_quasi-cyclic}, it is easy to prove that $\left\langle \mathbf{c_{1}}, \mathbf{c_{2}}\right\rangle_{s}=\mathbf{c_{1}} \cdot \Omega _{mn} \cdot \mathbf{c_{2}}^{T}=0$.
	Therefore, $\mathscr{C}$ and $\mathscr{C}^{\circ}$ are mutually symplectic orthogonal.
%	According to Definition \ref{one_quasi-cyclic_def}, $\mathscr{C}$ is a $1$-generator quasi-cyclic code has generator $(\left[g(x){f}_{0}(x)\right]$, $\left[g(x){f}_{1}(x)\right]$,$\ldots$, $\left[g(x){f}_{\ell -1}(x)\right])$, where $0\le j\le \ell -1$. 

 		%Since $\gcd(\bar{f}_{m-1}(x),x^n-1)=1$, the rank of $(F_{m-1}, F_{j})$ is $n$. In addition, 

 		By Lemma \ref{QCmatrix_spread}, the generator matrix of $\mathscr{C}^{\circ}$ also can be transformed into 
\begin{equation}
	G_1^{\perp_{s}}=
	\begin{pmatrix}
		F_{0}^{*}	& F_{1}^{*}\\
		A	& \mathbf{0}\\
		\mathbf{0}	& B\\
		\mathbf{0}&G_{g^{\perp_e}}
	\end{pmatrix},
\end{equation}
 		where $(F_{0}^{*},F_1^{*})$ denotes the first $(n-\deg(\gcd(x^n-1,\bar{f}_1(x)))-\deg(\gcd(x^n-1,\bar{f}_{0}(x))))$-row of $(F_{0},F_1)$, $A$ and $B$ are the generator matrices of cyclic codes generated by $\left[  \frac{x^n-1}{\gcd(x^n-1,\bar{f}_{1}(x))}\right]$ and $\left[ \frac{x^n-1}{\gcd(x^n-1,\bar{f}_{0}(x))}\right]$, respectively.
 		
 		According to Definition \ref{one_quasi-cyclic_def}, $\gcd(\bar{f}_{0}(x),\bar{f}_{1}(x),g^{\perp_e}(x))=1$. Given that $\gcd(f_0(x),g(x))=1$, we have $\gcd(\bar{f}_0(x),\tilde{g}(x) )=1$. As $\tilde{g}(x)g^{\perp_e}(x)=x^n-1$,  $\gcd(\bar{f}_{0}(x),\bar{f}_{1}(x),x^n-1)=1$ and $\operatorname{lcm}\left(g^{\perp_e}(x),\frac{x^n-1}{\gcd(x^n-1,\bar{f}_{0}(x))}\right) \equiv 0 \pmod{x^n-1}$ hold. This implies that, the rank of $(F_{0}, F_{1})$ is $n$ and the rows of $B$ and $G_{g^{\perp_e}}$ are linearly independent. Therefore, the rank of $G_1^{\perp_{s}}$ is $n+\deg(g(x))$. Thus, $\mathscr{C}$ and $\mathscr{C}^{\circ}$ are symplectic dual codes.
\end{proof}
%\begin{remark}\label{discuss_structure}
%	It is worth mentioning that Theorem \ref{Dual_Structure} only gives one form of the generator matrix of $\mathscr{C}^{\perp_s}$.
%	If we change the conditions in Theorem \ref{Dual_Structure}, the generator matrix of $\mathscr{C}^{\perp_s}$ can also be written in other different forms.
%	
%	 (1) For $\ell=2$, limiting $\gcd(\bar{f}_1(x),g(x))=1$, we can use a similar method to prove that $
%	\begin{pmatrix}
%		F_{0}	& F_{1}\\
%		G_{g^{\perp_e}}&
%	\end{pmatrix}_{n\times (n+\deg(g(x)))}$ is also a generator matrix of $\mathscr{C}^{\perp_s}$.
%
%	(2) For $\ell\ge4$, if $\gcd(\bar{f}_j(x),g(x))=1$, then $G_{g^{\perp_e}}$ can be written in other positions of the last line (the $j$-th block) of $G_2^{\perp_{s}}$.
%\end{remark}

%Obviously, according to Theorem \ref{Dual_Structure}, for $1$-generator quasi-cyclic codes, it is difficult to construct quasi-cyclic codes with large symplectic dual distances when the index $\ell \ge 4$. Therefore, we will only detail the $1$-generator quasi-cyclic codes with $\ell = 2$ as shown in the following theorem, which describes the upper and lower bounds on its minimum symplectic dual distance.

\begin{theorem}\label{SDQCBound} With the above notation,
	let $\mathscr{C}$ be a $1$-generator QC code with generator $(\left[f_0(x)g(x)\right],\left[f_1(x)g(x)\right])$ satisfying
 	$\gcd(f_0(x),g(x))=1$. 
  For $\alpha\in \mathbb{F}_q^*$, define $\mathcal{I}^{(\alpha)}_{\perp}(\mathscr{C})=\gcd(\bar{f}_1(x)+\alpha  \bar{f}_0(x),g^{\perp_{e}}(x))$. Let $\mathcal{S}_{\perp} =\{\mathcal{I}^{(\alpha)}_{\perp}(\mathscr{C}) \mid \mathcal{I}^{(\alpha)}_{\perp}(\mathscr{C})\ne 1, \forall\alpha\in\mathbb{F}_q^*\}$.
  %Let $\mathcal{S} =\{\mathcal{I}^{(\alpha)}_{\perp}(\mathscr{C})\mid \mathcal{I}^{(\alpha)}_{\perp}(\mathscr{C})=\gcd(\bar{f}_1(x)+\alpha  \bar{f}_0(x),g^{\perp_{e}}(x))\ne1, \alpha  \in \mathbb{F}_q^*\}$.
  
	Suppose that $\mathscr{C}_0$, $\mathscr{C}_1$, $\mathscr{C}_2$, $\mathscr{C}_3$, $\mathscr{C}_4$, $\mathscr{C}_5$, $\mathscr{C}_6$ and $\mathscr{C}_7$ are cyclic codes determined by 
 $\left[g^{\perp_{e}}(x)\right]$,
	$\left[\bar{f}_0(x)\right]$, $\left[\gcd(\bar{f}_1(x),g^{\perp_{e}}(x))\right]$,
	$\left[\frac{x^n-1}{\gcd( x^n-1,\bar{f}_1(x))}\right]$, 	$\left[\frac{x^n-1}{\gcd( x^n-1,\bar{f}_0(x))}\right]$, $\left[\gcd\left( \frac{x^n-1}{\gcd( x^n-1,\bar{f}_0(x))},g^{\perp_{e}}(x) \right)\right]$, $\left[ \operatorname{lcm}(\bar{f}_0(x),\gcd(\bar{f}_1(x),g^{\perp_{e}}(x))\right]$ and $\left[\bar{f}_0(x)\frac{g^{\perp_{e}}(x)}{\gcd(\bar{f}_1(x),g^{\perp_{e}}(x))}\right]$, respectively.

Define $\mathcal{D}_{\perp}(\mathscr{C})$ as in Equation (\ref{The.6.Eq.1}).
\begin{figure*}
	\begin{equation}\label{The.6.Eq.1}
		\mathcal{D}_{\perp}(\mathscr{C})=\max \left\{ \left\lceil\left(d_H(\mathscr{C}_1)+d_H(\mathscr{C}_2)+(q-|\mathcal{S}_{\perp} |-1)+\sum\limits_{\mathcal{I}^{(\alpha)}_{\perp}(\mathscr{C})\in \mathcal{S}_{\perp} } d_H\left(\left[ \mathcal{I}^{(\alpha)}_{\perp}(\mathscr{C}) \right]\right)\right)/ q\right\rceil, d_H(\mathscr{C}_1),d_H(\mathscr{C}_2)\right\}.
	\end{equation}
\end{figure*}
Writing $d_{\mathrm{upper}}=\min\left\{d_H(\mathscr{C}_0),d_H(\mathscr{C}_3),d_H(\mathscr{C}_4)\right\}$ and  $d_{\mathrm{lower}}$ as in Equation (\ref{Dual_lower_Bound}),
\begin{figure*}\begin{equation}\label{Dual_lower_Bound}
	d_{\mathrm{lower}}=
	\left\{
	\begin{array}{ll}
		\min \left\{ d_H(\mathscr{C}_3),d_H(\mathscr{C}_5),\mathcal{D}_{\perp}(\mathscr{C}) \right\},& \mathcal{S}_{\perp} =\emptyset, g(x)\mid \bar{f}_1(x),                                                           \\
		\min \left\{ d_H(\mathscr{C}_3),d_H(\mathscr{C}_5),d_H(\mathscr{C}_7),\mathcal{D}_{\perp}(\mathscr{C}) \right\},& \mathcal{S}_{\perp} =\emptyset, g(x)\nmid \bar{f}_1(x),                                       \\
		\min\left\{ d_H(\mathscr{C}_3),d_H(\mathscr{C}_5), d_H(\mathscr{C}_6), \mathcal{D}_{\perp}(\mathscr{C})\right\},&q=2,\mathcal{S}_{\perp} \ne\emptyset, g(x)\mid \bar{f}_1(x),                                   \\
		\min\left\{ d_H(\mathscr{C}_3),d_H(\mathscr{C}_5), d_H(\mathscr{C}_6),d_H(\mathscr{C}_7), \mathcal{D}_{\perp}(\mathscr{C})\right\},& q=2,\mathcal{S}_{\perp} \ne\emptyset,g(x)\nmid \bar{f}_1(x),\\
		\min \left\{ d_H(\mathscr{C}_3),d_H(\mathscr{C}_5),\max\{d_H(\mathscr{C}_1),d_H(\mathscr{C}_2)\} \right\},&q>2,\mathcal{S}_{\perp}\ne \emptyset,g(x)\mid \bar{f}_1(x),                     \\
		\min \left\{ d_H(\mathscr{C}_3),d_H(\mathscr{C}_5),d_H(\mathscr{C}_7),\max\{d_H(\mathscr{C}_1),d_H(\mathscr{C}_2)\} \right\},&q>2,\mathcal{S}_{\perp}\ne \emptyset,g(x)\nmid \bar{f}_1(x). \\
	\end{array}
	\right.
\end{equation} \end{figure*}
then we have $d_{\mathrm{lower}}\le d_s(\mathscr{C}^{\perp_s}) \le d_{\mathrm{upper}}$.
\end{theorem}
\begin{proof}
Let $t=(n-\deg(\gcd(x^n-1,\bar{f}_1(x)))-\deg(\gcd(x^n-1,\bar{f}_{0}(x))))$ and $\mathscr{C}_8$ be the cyclic code generated by $\left[ \bar{f}_1(x) \right]$. 
	Since $\gcd(f_0(x),g(x))=1$,
	by Theorem \ref{Dual_Structure}, the symplectic dual of $\mathscr{C}$ is a 2-generator QC code with generator matrix
	\begin{equation}
		G^{\perp_{s}}
  % =\left(\begin{array}{cc}
		% 	F_0 & F_1 \\
		% 	\mathbf{0} & G_{g^{\perp_e}}
		% \end{array}\right)
  =
  \begin{pmatrix}
		F_{0}^{*}	& F_{1}^{*}\\
		A	& \mathbf{0}\\
		\mathbf{0}	& B\\
		\mathbf{0}&G_{g^{\perp_e}}
	\end{pmatrix},
	\end{equation}
where
$(F_{0}^{*},F_1^{*})=\mathrm{Circ}_t\left( \left[\bar{f}_0(x)\right],\left[\bar{f}_1(x)\right] \right)$, $A$ and $B$ are the generator matrices of cyclic codes generated by $\left[  \frac{x^n-1}{\gcd(x^n-1,\bar{f}_{1}(x))}\right]$ and $\left[ \frac{x^n-1}{\gcd(x^n-1,\bar{f}_{0}(x))}\right]$, and $G_{g^{\perp_e}}$ is the generator matrix of cyclic code generated by $[g^{\perp_e}(x)]$, respectively.  
	% where $(F_{0},F_1)$ and $(\mathbf{0},G_{g^{\perp_e}})$ are generator matrices of quasi-cyclic codes generated by $(\left[ \bar{f}_{0}(x)\right], \left[ \bar{f}_{1}(x)\right])$ and $(\left[ 0\right], \left[ G_{g^{\perp_e}}\right])$, respectively. 
Accordingly, any codeword of $\mathscr{C}^{\perp_{s}}$ can be expressed as $\mathbf{c}=(\left[a(x)\bar{f}_0(x)\right],\left[a(x)\bar{f}_1(x)+b(x)g^{\perp_{e}}(x)\right])$, where $a(x),b(x)\in\mathbb{R}_{q,n}$. 
	By \cite{ling2010generalization}, we have
 % \begin{equation}
     $q\mathbf{w}_{s}(\mathbf{c})=\mathbf{w}_{H}(\left[a(x)\bar{f}_0(x)\right])+\mathbf{w}_{H}(\left[a(x)\bar{f}_1(x)+b(x)g^{\perp_{e}}(x)\right])+\sum\limits_{\alpha  \in \mathbb{F}_q^*} \mathbf{w}_{H}(\left[a(x)(\bar{f}_1(x)+\alpha   \bar{f}_0(x))+b(x)g^{\perp_{e}}(x)\right]).$
 % \end{equation}
	
	
	Similar to the proof of Theorem \ref{SQCBound}, if $a(x)\bar{f}_1(x)\equiv  0 \pmod{x^n-1}$ or $a(x)\bar{f}_0(x) \equiv  0 \pmod{x^n-1}$, then $\mathbf{c}=(\left[a(x)\bar{f}_0(x)\right],\left[b(x)g^{\perp_{e}}(x)\right])$ or $(\mathbf{0}_n,\left[a(x)\bar{f}_1(x)+b(x)g^{\perp_{e}}(x)\right])$.
	At this point, we have $\min \left\{ d_H(\mathscr{C}_3),d_H(\mathscr{C}_5) \right\}  \le \mathbf{w}_{s}(\mathbf{c}) \le 
	\min\left\{d_H(\mathscr{C}_0),d_H(\mathscr{C}_3),d_H(\mathscr{C}_4)\right\}$.

	
	%Let $\mathscr{C}^{\circ}$ be a $1$-generator quasi-cyclic code determined by $(\left[\bar{f}_0(x)\right],\left[\gcd(\bar{f}_1(x),g^{\perp_{e}}(x))\right])$. 
  
 If $a(x)\bar{f}_0(x)\not\equiv 0 \pmod{x^n-1}$ and $a(x)\bar{f}_1(x) \not\equiv 0 \pmod{x^n-1}$, there are the following scenarios to discuss. 

	Firstly,  if $g(x)\mid \bar{f}_1(x)$, then $\operatorname{lcm}(\bar{f}_1(x),g^{\perp_e}(x))\equiv 0 \pmod{x^n-1}$. 
 It implies that $\mathscr{C}_0\cap \mathscr{C}_8=\{ \mathbf{0}_{n}\}$. 
 %, so $\left[ g^{\perp_e}(x) \right]\cap\left[ \bar{f}_1(x) \right]=\{\mathbf{0}\}$.
	 If $g(x)\nmid \bar{f}_1(x)$, then the two cyclic codes $\mathscr{C}_0$ and $\mathscr{C}_8$ have nontrivial intersection.
 Therefore, codewords of the form $(\left[a(x)\bar{f}_0(x)\right],\mathbf{0}_n)$ may also exist in $\mathscr{C}^{\perp_s}$. 
 In this case, there exists $\frac{g^{\perp_{e}}(x)}{\gcd(\bar{f}_1(x),g^{\perp_{e}}(x))} \mid a(x)$. Thus, when $g(x)\nmid \bar{f}_1(x)$, we can only deduce $d_s(\mathscr{C}^{\perp_s})\ge d_H(\mathscr{C}_7)$. 
	
	Secondly, if for all $ \alpha  $ in $\mathbb{F}_q^*$, $\gcd(\bar{f}_1(x)+\alpha   \bar{f}_0(x),g^{\perp_{e}}(x))= 1$, i.e., $\mathcal{S}_{\perp} =\emptyset$, then 
	\begin{equation}\label{EDSBound1}
		d_s(\mathscr{C}^{\perp_s})\ge \left(d_H(\mathscr{C}_1)+d_H(\mathscr{C}_2)+(q-1) \right)/ q.
	\end{equation}
	
	Thirdly, if there exists $\alpha  $ in $\mathbb{F}_q^*$ such that  $\gcd(\bar{f}_1(x)+\alpha   \bar{f}_0(x),g^{\perp_{e}}(x))\ne 1$, i.e., $\mathcal{S}_{\perp} \ne\emptyset$, then there may exist $\mathrm{Supp}(\left[a(x)\bar{f}_0(x)\right]) =\mathrm{Supp}(\left[a(x)\bar{f}_1(x)+b(x)g^{\perp_{e}}(x)\right])$.
	Thus, for $q>2$, we can only deduce $d_s(\mathscr{C}^{\perp_s})\ge \max\left\{ d_H(\mathscr{C}_1),d_H(\mathscr{C}_2) \right\}$.
	But for $q=2$, we have
	$d_s(\mathscr{C}^{\perp_s})\ge d_H( \left[ \operatorname{lcm}(\bar{f}_0(x),\gcd(\bar{f}_1(x),g^{\perp_{e}}(x))) \right])$. 
	
	For the case of $\mathrm{Supp}(\left[a(x)\bar{f}_0(x)\right])\ne$ 
	$ \mathrm{Supp}(\left[a(x)\bar{f}_1(x)+\right. $ $\left.b(x)g^{\perp_{e}}(x)\right])$,
	we have Equation (\ref{EDSBound2}) holds.
	

	
	 \begin{figure*}
	\begin{equation}\label{EDSBound2}
		d_s(\mathscr{C}^{\perp_s})\ge \left(d_H(\mathscr{C}_1)+d_H(\mathscr{C}_2)+(q-|\mathcal{S}_{\perp} |-1)+\sum\limits_{\mathcal{I}^{(\alpha)}_{\perp}(\mathscr{C})\in \mathcal{S}_{\perp} } d_H\left(\left[ \mathcal{I}^{(\alpha)}_{\perp}(\mathscr{C}) \right]\right)\right)/ q.
	\end{equation} \end{figure*} 
	
	Moreover, at this time, $d_s(\mathscr{C}^{\perp_s})$ must be greater than or equal to $\max\left\{ d_H(\mathscr{C}_1),d_H(\mathscr{C}_2) \right\}$ when Equations (\ref{EDSBound1}) and (\ref{EDSBound2}) are satisfied.
 Therefore, we have $d_s(\mathscr{C})\ge\mathcal{D}_{\perp}(\mathscr{C})$.
 
	From the combination of the above conditions, one can obtain upper and lower bounds for the minimum symplectic distance of $\mathscr{C}^{\perp_s}$. This concludes the proof.
\end{proof}
%
%\begin{remark}
%	In 2023, Dastbasteh et al. also studied a lower bound on the minimum distance for a class of additive cyclic codes (see Theorem 3.9 in \cite{dastbasteh2023polynomial}), which is similar to our Theorem \ref{SDQCBound} for quasi-cyclic codes  ($\Phi(\mathscr{C}^{\perp_s} )$ is also additive cyclic codes).
%	%Recently, Dastbasteh et al. also studied the minimum distance lower bound on a class of additive cyclic codes (see Theorem 3.9 in \cite{dastbasteh2023polynomial}), which is similar to our Theorem \ref{SDQCBound} for $\Phi(\mathscr{C}^{\perp_s})$ is additive cyclic codes. 
%	However, we find that Theorem 3.9 in \cite{dastbasteh2023polynomial} seems to miss some cases. For example, let $q = 2$ and $n = 21$, set $g(x)=x^{20} + x^{19} + x^{18} + x^{17} + x^{15}$, $k(x)=x^{20} + x^{16} + 1$ and $h(x)=x^9 + x^8 + x^7 + x^2 + x + 1$ (the symbol represents the same as Theorem 3.9 in \cite{dastbasteh2023polynomial}). Since $g(x)\nmid \bar{f}_1(x)$ has not been taken into account in Theorem 3.9 in \cite{dastbasteh2023polynomial}, it gives a lower bound on minimum distance of additive cyclic code $\left[ g(x)+wk(x),wh(x) \right]$, denoted by $\mathscr{C}_a$ , is $3$, but the actual distance is $2$, which is a contradiction.
%	According to Theorem \ref{SDQCBound}, the lower bound on the minimum symplectic distance of $\Phi^{-1}(\mathscr{C}_a)$ is $2$ (by setting $\bar{f}_0(x)=g(x),\bar{f}_1(x)=k(x)$ and $g^{\perp_{e}}(x)=h(x)$). 
%	%Therefore, our bound also can be seen as an improvement and correction of Theorem 3.9 in \cite{dastbasteh2023polynomial}.
%\end{remark}


%In Remark \ref{remark2}, we have give a method to construct specific $f(x)$, which satisfy $f(x)=\bar{f}(x)$. Thus, in the process of constructing symplectic self-orthognal codes via Theorem \ref{theorem_(1,2)-QC}, if $g^{\perp_{e}}(x)\nmid g(x)$, one can set $f_{i}(x)=\bar{f}_{i}(x)$, where $i =0$ or $1$, which will makes $g^{\perp_{e}}(x) \mid g(x)\left(f_{0}(x) \bar{f}_{1}(x)-f_{1}(x) \bar{f}_{0}(x)\right)$ holds constantly. This will significantly reduce the computational power and time required to search for symplectic self-orthogonal codes with good parameters.

%In addition, by means of Lemma \ref{CRSS}, we have the following theorem to help deduce QECCs from symplectic self-orthogonal quasi-cyclic codes.


%\begin{theorem}\label{QC-QECC}
% Let $\mathscr{C}$ be $1$-generator quasi-cyclic code in Definition  \ref{(1,2)-QC} and satisfy Theorem \ref{theorem_(1,2)-QC}, then there exist a QECC with parameters  $\left[\left[n,deg(g(x)),{d}_{s}\right]\right]_q$, where ${d_{s}}=\min \{swt(c)\mid c\in{{\mathscr{C}}^{{\bot}_{s}}}\backslash \mathscr{C}\}$. 
%\end{theorem}
%	
%By Theorem \ref{one_quasi-cyclic}, the necessary and sufficient condition for $\mathscr{C}$ to be symplectic self-orthogonal is
%$g^{\perp_{e}}(x) \mid g(x)\left(f_{0}(x) \bar{f}_{1}(x)-f_{1}(x) \bar{f}_{0}(x)\right)$. 

\begin{remark}\label{improve_bound_Galindo}
Theorem \ref{SDQCBound} estimates the range of the minimum symplectic distance for $2$-generator QC codes $\mathscr{C}^{\perp_s}$ in Theorem \ref{Dual_Structure}. 
It should be noted that Galindo et al. also studied the lower bound on minimum symplectic distance for a class of  $2$-generator QC codes with generators $(\left[g_1(x)\right],\left[f(x)g_1(x)\right])$ and $([0],\left[g_2(x)\right])$, where $g_i(x) \mid (x^n-1)$ and $f(x)\in \mathbb{R}_{q,n}$ (please see Proposition 4 in \cite{galindo2018quasi}). 
It is easy to verify that if we constraint $g_1(x)=1$, then Galindo's codes are a subclass of $\mathscr{C}^{\perp_s}$ in Theorem \ref{Dual_Structure}.   Therefore,
Theorem \ref{SDQCBound} can be seen as a valid  generalization of Proposition 4 in \cite{galindo2018quasi}. 
\end{remark}

\begin{example}\label{good_Ex2}
	Let $q=2$ and $n=15$. Set $g(x)=x^4 + x + 1$, $f_0(x)=x^{13} + x^{12} + x^{11} + x^8 + x^7 + x^4 + x^3 + x^2 + 1$ and $f_1(x)=x^{13} + x^9 + x^8 + x^7 + x^6 + x^2 + 1$, then $g^{\perp_e}(x)=x^{11} + x^{10} + x^9 + x^8 + x^6 + x^4 + x^3 + 1$. It is easy to check that $\gcd(f_0(x),f_1(x))=1$, $\gcd(\bar{f}_0(x),g(x))=1$ and 
 $g^{\perp_e}(x) \mid g(x)(f_{0}(x) \bar{f}_{1}(x)-f_{1}(x) \bar{f}_{0}(x))$, so $(\left[g(x)f_0(x)\right],\left[g(x)f_1(x)\right])$ generates a symplectic self-orthogonal $\left[30,11\right]_2$ QC code $\mathscr{C}$.
 

Since $\mathcal{S}_{\perp} =\{\mathcal{I}^{(1)}_{\perp}(\mathscr{C})=\gcd(\bar{f}_1(x)+ \bar{f}_0(x),g^{\perp_{e}}(x))= x + 1\}$ and $\operatorname{lcm}(\bar{f}_1(x),g^{\perp_e}(x))\equiv 0 \pmod{x^n-1}$, we get


 	\begin{equation} \label{improve_boundD}
 			\begin{array}{rl}
 	\min\left\{ d_H(\mathscr{C}_3),d_H(\mathscr{C}_5), d_H(\mathscr{C}_6), \mathcal{D}_{\perp}(\mathscr{C})\right\}  \le d_s(\mathscr{C}^{\perp_s})\\
 	 \le 
 	\min\left\{d_H(\mathscr{C}_0),d_H(\mathscr{C}_3),d_H(\mathscr{C}_4)\right\},
 \end{array}
 \end{equation}
where 
 $
 	\mathcal{D}_{\perp}(\mathscr{C})=\max \left\{ \left\lceil\left(d_H(\mathscr{C}_1)+d_H(\mathscr{C}_2)+ d_H\left(\left[ \mathcal{I}^{(1)}_{\perp}(\mathscr{C}) \right]\right)\right)/ 2\right\rceil\right.
 	,$ $\left.  d_H(\mathscr{C}_1),d_H(\mathscr{C}_2)\right\}.
$
% and $\mathscr{C}_0$, $\mathscr{C}_1$, $\mathscr{C}_2$, $\mathscr{C}_3$, $\mathscr{C}_4$, $\mathscr{C}_5$ and $\mathscr{C}_6$ are cyclic codes determined by $\left[g^{\perp_{e}}(x)\right]$,
%  $\left[\bar{f}_0(x)\right]$, $\left[\gcd(\bar{f}_1(x),g^{\perp_{e}}(x))\right]$,
%  $\left[\frac{x^{15}-1}{\gcd( x^{15}-1,\bar{f}_1(x))}\right]$, 	$\left[\frac{x^{15}-1}{\gcd( x^{15}-1,\bar{f}_0(x))}\right]$, $\left[\gcd\left( \frac{x^{15}-1}{\gcd( x^{15}-1,\bar{f}_0(x))},g^{\perp_{e}}(x) \right)\right]$ and $\left[\operatorname{lcm}(\bar{f}_0(x),\gcd(\bar{f}_1(x),g^{\perp_{e}}(x))\right]$, respectively.
Here, we list the generator polynomials and parameters of the cyclic codes involved.
 \begin{itemize}
      \item  $\mathscr{C}_0$: Generator polynomial $\left[g^{\perp_{e}}(x)\right]$. 
      Parameters $\left[15,4,8\right]_2$.
      \item  $\mathscr{C}_1$: Generator polynomial $\left[\bar{f}_0(x)\right]=\left[x^6 + x^4 + x^3 + x^2 + 1\right]$. Parameters $\left[15,9,4\right]_2$.
      \item  $\mathscr{C}_2$: Generator polynomial $\left[\gcd(\bar{f}_1(x),g^{\perp_{e}}(x))\right]=\left[x^4 + x + 1\right]$. Parameters $\left[15,11,3\right]_2$.
     \item $\mathscr{C}_3$: Generator polynomial $\left[\frac{x^{15}-1}{\gcd( x^{15}-1,\bar{f}_1(x))}\right]=\left[x^7 + x^6 + x^5 + x^2 + x + 1\right]$. Parameters $\left[15,8,4\right]_2$.
      \item $\mathscr{C}_4$: Generator polynomial $\left[\frac{x^{15}-1}{\gcd( x^{15}-1,\bar{f}_0(x))}\right]=\left[x^9 + x^7 + x^6 + x^3 + x^2 + 1\right]$. Parameters $\left[15,6,6\right]_2$.
      \item  $\mathscr{C}_5$: Generator polynomial $\left[\gcd\left( \frac{x^{15}-1}{\gcd( x^{15}-1,\bar{f}_0(x))},g^{\perp_{e}}(x) \right)\right]$ $=\left[x^5 + x^4 + x^2 + 1\right]$. Parameters $\left[15,10,4\right]_2$.
      \item  $\mathscr{C}_6$: Generator polynomial $\left[\operatorname{lcm}(\bar{f}_0(x),\gcd(\bar{f}_1(x),g^{\perp_{e}}(x))\right]=\left[x^{10} + x^8 + x^5 + x^4 + x^2 + x + 1\right]$. Parameters $\left[15,5,7\right]_2$.
      \item  Generator polynomial $\left[\mathcal{I}^{(1)}_{\perp}(\mathscr{C})\right]= \left[x + 1\right]$. Parameters $\left[15,14,2\right]_2$.
 \end{itemize}

 
 It follows directly that the lower and upper bounds on the minimum symplectic distance of $\mathscr{C}^{\perp_s}$ are both $4$ based on the parameters of the relevant cyclic code, so the parameters of $\mathscr{C}^{\perp_s}$ are $\left[30,19,4\right]_2^s$. By using Magma \cite{bosma1997magma}, we can compute that the specific parameters of $\mathscr{C}^{\perp_s}$ are indeed $\left[30,19,4\right]_2^s$.
Thus, by Theorem \ref{CRSS}, we can obtain a binary QECC with parameters $\left[\left[15,4,4\right]\right]_2$, which is an optimal QECC according to \cite{Grassltable}.
\end{example}



% \begin{remark}
% 	Theorems \ref{one_quasi-cyclic} and \ref{theo-h-QC} can help us to determine which quasi-cyclic codes are symplectic self-orthogonal, thus greatly reducing the search space. Furthermore, Theorems \ref{SQCBound} and \ref{SDQCBound} can help to determine the range of the minimum symplectic distances of $1$-generator quasi-cyclic codes with index two and their symplectic dual codes (the distances of cyclic codes can be determined by previous bounds in \cite{huffman2010fundamentals}), with lower cost, thereby further helping us to reduce the search space. 
% 	Moreover, although we do not derive the symplectic dual codes of $h$-quasi-cyclic codes and related bounds, we can still utilize them to efficiently construct symplectic self-orthogonal codes with the help of Theorem \ref{theo-h-QC}.  
% \end{remark}

By utilizing the results above, we devise many superior symplectic self-orthogonal codes via computerized search, corresponding to 117 record-breaking QECCs, see Table \ref{all_QECCs}.
Specifically, we give the specific generators of related codes and their corresponding parameters in Tables \ref{F_QECCs} and \ref{T_QECCs}. 
In order to save space, we give them in abbreviated form and represent the coefficient polynomials in ascending order, where the indexes of elements represent consecutive elements of the same number. For example, polynomial $1 + x^2 + x^3 + x^4$ over $\mathbb{F}_2$ is denoted as $101^3$. The specific parameters of corresponding codes are calculated by using Magma \cite{bosma1997magma}.

%In Table 2 and their derived codes (by Lemma \ref{propagation_rules}) are given in Table 2, all of which improve the lower bound on the minimum distances in code table \cite{Grassltable}.


\begin{remark}
	The new binary QECCs presented in this paper have been updated in Grassl's code table \cite{Grassltable}.  Additionally, after we completed this paper, Grassl et al. \cite{ezerman2024characterization} systematically constructed a large number of good QECCs using QC codes, and extended the quantum code table to $q\le 8$ and $n\le 100$.
\end{remark}


%\begin{remark}
%	For example, with $q=3$ and $n=13$, it is easy to obtain ternary optimal QECCs with parameters $\left[\left[13,0,6\right]\right]_3$, $\left[\left[13,3,5\right]\right]_3$, $\left[\left[13,6,4\right]\right]_3$, and $\left[\left[13,8,3\right]\right]_3$, which all attain the (upper) bounds on the minimum distance in \cite{Grassltable}.
%	This is generated as follows, with the polynomials in the same abbreviated form as in Tables \ref{F_QECCs} and \ref{T_QECCs} to save space, and $\mathbb{F}_3=\{0,1,2\}$.
%	\begin{itemize}
%		\item $\left[\left[13,0,6\right]\right]_3$: $g(x)=1$, $f_0(x)=0212^{2}10^{2}12^{2}12$, 
%		$f_1(x)=02^{2}1^{2}0^{4}1^{2}2^2$.
%		\item $\left[\left[13,3,5\right]\right]_3$: $g(x)=201^2$, $f_0(x)=2^{4}0^{2}1^{2}0^{2}2^3$, 
%		$f_1(x)=2^{2}02^{2}10^{2}12^{2}02$.
%		\item $\left[\left[13,6,4\right]\right]_3$: $g(x)=1^{2}0^{2}101$, $f_0(x)=12^{3}0^{2}2^{2}0^{2}2^3$, 
%		$f_1(x)=102020^{4}202$.
%		\item $\left[\left[13,8,3\right]\right]_3$: $g(x)=1012^{2}0201$, 
%		$f_0(x)=10212^{6}12$, 
%		$f_1(x)=1^{3}0^{2}21^{2}20^{2}1^2$.
%	\end{itemize}
%	\begin{table*}\left[!t\right]
%		\centering
%		\begin{threeparttable}
%			\caption{One-generator symplectic self-orthogonal quasi-cyclic codes\tnote{4} with $\ell=2$  and Ternary  QECCs\tnote{5}}\label{Ternary_QECCs}
%			\centering
%			\footnotesize
%			\begin{tabular}{cccc}
%				\hline
%				No. &     QECCs      & Theorem \ref{SDQCBound} &                                         $g(x),f_0(x),f_1(x)$                                          \\ \hline
%				 1  & $\left[\left[13,0,6\right]\right]_3$ &     $5\le d_s\le 6$     & \makecell\left[c\right]{$1$, $2^{1}0^{3}1^{1}2^{1}1^{2}2^{1}1$, $2^{3}1^{1}0^{1}2^{1}0^{2}2^{1}0^{1}1^{1}2^2$} \\
%				 2  & $\left[\left[13,3,5\right]\right]_3$ &    $\left[84,49,10\right]_2^s$     &                                      \makecell\left[c\right]{$10^{7}101$, }                                      \\
%				 3  & $\left[\left[13,6,4\right]\right]_3$ &   $\left[ 90, 66, 7 \right]_2^s$   &                         \makecell\left[c\right]{$1^{4}0101^{4}0^{2}101^{2}0^{2}10^{3}1$}                         \\
%				 4  & $\left[\left[13,8,3\right]\right]_3$ &    $\left[110,75,9\right]_2^s$     &               \makecell\left[c\right]{$10^{2}101^{2}010^{2}10^{3}1^{2}0^{4}10^{2}10^{2}1010101^2$}               \\ \hline
%			\end{tabular}
%			\begin{tablenotes}    %这行要添加， 从这开始
%				\footnotesize               %这行要添加
%				\item\left[4\right] These $2$-generator quasi-cyclic codes with generator elements $(\left[g_1(x)f(x)\right],\left[g_1(x)\right])$ and $(\left[g_2(x)\right],\left[g_2(x)f(x)\right])$.        %这行要添加
%				\item\left[5\right] $\left[\left[55,20,9\right]\right]_2$ and $\left[\left[63,31,8\right]\right]_2$ were presented in our first version of \cite{guan2021new}, but were not formally published in \cite{guan2022new}; here we formally submit them for publication.  In addition, during our revision of this paper \cite{guan2022symplectic}, we noticed that Padmapani Seneviratne also obtained $\left[\left[40,6,10\right]\right]_2$ and updated it in the code table \cite{Grassltable}.        %这行要添加
%			\end{tablenotes}            %这行要添加
%		\end{threeparttable}       %这行要添加，到这里结束
%	\end{table*}
%\end{remark}



%\begin{definition}\label{(2,2)-QC}
%	Suppose $g_1(x)$, $g_2(x)$, $f(x)$ are polynomials in $\mathbb{R}_{q,n}$, where $g_i(x) \mid\left(x^{n}-1\right)$, $1 \leq i \leq 2$.  Let $\mathscr{C}$ be a 2-generator quasi-cyclic code of index two with generators
%	$\left(\left[g(x) f(x)\right],\left[g(x)\right]\right)$, $\left(\left[g(x)\right],\left[g(x) f(x)\right]\right)$.  
%	Then, generator matrix $G$ of $\mathscr{C}$ has the following form:
%	\begin{equation}
%		G=\begin{pmatrix}  
%			G^{\prime }_{1} & G_1 \\  
%			G_{2} & G^{\prime }_{2} 
%		\end{pmatrix} ,
%	\end{equation}
%	where $G_{i}$ and $G^{\prime}_{i}$ are $n \times n$ circulant matrices generated by $\left[g_i(x)\right]$ and $\left[g_{i}(x) f(x)\right]$, respectively.
%\end{definition}
%
%
%	Let $\mathscr{C}$ be $2$-generator quasi-cyclic codes in Definition \ref{(2,2)-QC}. 
%	It follows from Theorem \ref{theo-h-QC} that necessary and sufficient conditions for $\mathscr{C}$ to be a symplectic self-orthogonal code are


%//-----------------------------------------------------
%\begin{corollary}\label{theorem_(2,2)-QC}
%	Let the symbols be the same as above and let $\mathscr{C}$ be a $2$-generator quasi-cyclic code with index $2$ in Definition 2. If it satisfies
%	\begin{equation}\label{eq_2h}
%		\left\{
%		\begin{array}{c} 
%			g_i^{\perp_e}(x) \mid g_i(x)(f(x) - \bar{f}(x)), \\ 
%			g_2^{\perp_e}(x) \mid g_1(x)(f(x)\bar{f}(x)-1).
%		\end{array}
%		\right.
%	\end{equation}
%	then $\mathscr{C}$ is a symplectic self-orthogonal code.
%\end{corollary}
%\begin{proof}
%	By Definition 2 and Theorem \ref{theo-h-QC}, this corollary holds.
%	%	In addition, because of $\gcd(f_{j}(x),\frac{x^n-1}{g(x)} )=1$, 
%	%	so $g(x)=\gcd(g(x) f_{0}(x), g(x) f_{1}(x)$, $x^{n}-1)$. 
%	%	Therefore, the dimension of $\mathscr{C}$ is $n-\deg(g(x))$.
%\end{proof}
%Obviously, by choosing the appropriate $g_i(x)$ and $f(x)$, it is possible to construct the symplectic self-orthogonal codes very efficiently. 
%Thus, we employ quasi-cyclic codes in Definition \ref{(2,2)-QC} to construct five good symplectic self-orthogonal codes in Table \ref{T_QECCs}, all of which correspond to record-breaking QECCs. 
%Also, as in Table \ref{F_QECCs}, generator polynomials in Table \ref{T_QECCs} are given in abbreviated form.
%An example of the construction process is given below. 
%\begin{example}
%	Set $q=2$, $n=42$. Let $g_1(x)=x^{10} + x^8 + 1$,
%	$g_2(x)=x^{39} + x^{38} + x^{37} + x^{35} + x^{32} + x^{31} + x^{30} + x^{28} + x^{25} + x^{24} + x^{23} + 
%	x^{21} + x^{18} + x^{17} + x^{16} + x^{14} + x^{11} + x^{10} + x^9 + x^7 + x^4 + x^3 + x^2 + 
%	1$, choose
%	$f(x)=x^{41} + x^{40} + x^{31} + x^{30} + x^{29} + x^{24} + x^{23} + x^{22} + x^{21} + x^{20} + x^{19} + x^{18} + x^{13} + x^{12} + x^{11} + x^2 + x + 1$. 
%	It is easy to check that $g_1(x)$, $g_2(x)$, $f(x)$ satisfy equation (\ref{eq_2h}), so we can get an $\left[84,35\right]_2^s$ symplectic self-orthogonal code, whose symplectic dual's weight distribution is  
%	$w(z)=1+2667z^{10}+23016z^{11}+173376z^{12}+1172808z^{13}+7432512z^{14}+\dots+3209199603z^{42}$.
%	By Theorem \ref{CRSS}, a binary QECC with parameters  $\left[\left[42,7,10\right]\right]_2$  can be provided.
%	
%	Observe that a code with parameters  $\left[\left[42,7,9\right]\right]_2$ is the best-known binary QECC with length $42$ and dimension $7$ in \cite{Grassltable}. Therefore, the corresponding minimum distance record can be improved to $10$.
%\end{example}
%
%Finally, we give, in Table \ref{all_QECCs}, all the record-breaking QECCs obtained in this paper, some of which are derived codes (via Lemma \ref{propagation_rules}).
%//-----------------------------------------------------
\section{Conclusion}\label{sec5}

In this work, we established the necessary and sufficient conditions for QC codes with even 
index to be symplectic self-orthogonal. 
Furthermore, through the decomposition of the code space of QC codes, we established lower and upper bounds on the minimum symplectic distances of a class of $1$-generator QC codes with index two and their symplectic dual codes. 
This contributed to improving the results presented in \cite{galindo2018quasi}, \cite{dastbasteh2023polynomial}, and \cite{GuanAdditive2023}, see Remarks \ref{improve_bound} and \ref{improve_bound_Galindo}, and Examples \ref{good_Ex1} and \ref{good_Ex_for additive}.
As an application, we constructed 117 record-breaking binary QECCs from symplectic self-orthogonal QC codes.
%It should be supplemented that our theory is applicable to any field, but since the main purpose of this paper is to improve the records in \cite{Grassltable}, our results are all binary QECCs.
%The theory of this paper also performs well for non-binary QECCs. 

%Quasi-cyclic codes have a rich algebraic structure, and based on our computational results (see Tables \ref{T_QECCs} and \cite{guan2022new}), it is also valuable to study the construction of QECCs using multi-generator quasi-cyclic codes. 
%In addition, cyclic codes \cite{huffman2010fundamentals} and quasi-cyclic codes \cite{zeh2014decoding} are both enjoy efficient encoding and decoding algorithms, and simple additive cyclic QECC $[[7,1,3]]_2$ (\cite{abu2014relative,tansuwannont2021fault}) also can be decoded and applied to fault-tolerant quantum error-correction. 
%Therefore, it is interesting to study how to decode additive cyclic QECCs in this paper efficiently.

%In addition, we present a practical approach for constructing symplectic self-orthogonal codes.
%%Furthermore, we propose a class of symplectic self-orthogonal quasi-cyclic codes that improve the additive construction of Calderbank, Theorem 14 in \cite{calderbank1998quantum}.
%Finally, more than 100 record-breaking QECCs are constructed, demonstrating our method's effectiveness.

%Cyclic codes \cite{huffman2010fundamentals} and quasi-cyclic codes \cite{zeh2014decoding} are both enjoy efficient encoding and decoding algorithms, and as a simple additive cyclic QECC $[[7,1,3]]_2$ also can be decoded    
%
%In the future, it is more challenging to derive lower and upper bounds on the minimum symplectic distance for $h$-generator quasi-cyclic codes ($h\ge 2$), but according to our computational results (see Table \ref{T_QECCs} and \cite{guan2022new}), $h$-generator quasi-cyclic codes still have significant research value. 
%
%Quasi-cyclic codes have a rich algebraic structure. Consequently, it is possible to construct symplectic self-orthogonal codes with more complex algebraic structures, which may improve the records in code table \cite{Grassltable}.
%Therefore, it will be interesting to determine the lower and upper bounds of the symplectic distance of $h$-generator quasi-cyclic codes.
%Moreover, from Examples \ref{good_Ex1} and \ref{good_Ex2}, it is possible to derive quasi-cyclic codes with good symplectic distances from Theorems \ref{SQCBound} and \ref{SDQCBound} by using special cyclic codes. Therefore, it will also be interesting to study the use of special cyclic codes to construct quasi-cyclic codes with considerable lengths and symplectic distances.






\begin{table*}[!t]
	\centering
	\begin{threeparttable}
		\caption{ New binary QECCs  and derived record-breaking codes\tnote{1}$^{,}$\tnote{2}}\label{all_QECCs}
		\footnotesize
		\begin{tabular}{cccccc}
			\hline
			NO. & Record-breaking  QECCs  & Code Table \cite{Grassltable} & NO. &  Record-breaking QECCs  & Code Table \cite{Grassltable} \\ \hline
			 1  & $[[40,6,10]]_2^{\star}$ &       $[[40,6,8]]_2^*$        & 60  & $[[63,31,8]]_2^{\star}$ &               -               \\
			 2  &     $[[41,6,10]]_2$     &        $[[41,6,8]]_2$         & 61  &   $[[73,18,13]]_2^*$    &       $[[73,18,12]]_2$        \\
			 3  &     $[[39,7,9]]_2$      &        $[[39,7,8]]_2$         & 62  &    $[[73,17,13]]_2$     &       $[[73,17,12]]_2$        \\
			 4  &     $[[39,6,9]]_2$      &        $[[39,6,8]]_2$         & 63  &    $[[73,16,13]]_2$     &       $[[73,16,12]]_2$        \\
			 5  &     $[[40,7,9]]_2$      &        $[[40,7,8]]_2$         & 64  &    $[[73,15,13]]_2$     &       $[[73,15,12]]_2$        \\
			 6  & $[[42,7,10]]_2^{\star}$ &        $[[42,7,9]]_2$         & 65  &    $[[73,14,13]]_2$     &       $[[73,14,12]]_2$        \\
			 7  &     $[[43,7,10]]_2$     &        $[[43,7,9]]_2$         & 66  &    $[[74,18,13]]_2$     &       $[[74,18,12]]_2$        \\
			 8  &     $[[44,7,10]]_2$     &        $[[44,7,9]]_2$         & 67  &    $[[74,17,13]]_2$     &       $[[74,17,12]]_2$        \\
			 9  &     $[[45,7,10]]_2$     &        $[[45,7,9]]_2$         & 68  &    $[[74,16,13]]_2$     &       $[[74,16,12]]_2$        \\
			10  &     $[[46,7,10]]_2$     &        $[[46,7,9]]_2$         & 69  &    $[[74,15,13]]_2$     &       $[[74,15,12]]_2$        \\
			11  &     $[[41,8,9]]_2$      &        $[[41,8,8]]_2$         & 70  &    $[[75,18,13]]_2$     &       $[[75,18,12]]_2$        \\
			12  &    $[[40,5,10]]_2^*$    &        $[[40,5,9]]_2$         & 71  &    $[[75,17,13]]_2$     &       $[[75,17,12]]_2$        \\
			13  &     $[[41,5,10]]_2$     &        $[[41,5,9]]_2$         & 72  &    $[[75,16,13]]_2$     &       $[[75,16,12]]_2$        \\
			14  &     $[[42,5,10]]_2$     &        $[[42,5,9]]_3$         & 73  &    $[[75,15,13]]_2$     &       $[[75,15,12]]_2$        \\
			15  &     $[[39,6,9]]_2$      &        $[[39,6,8]]_4$         & 74  &    $[[76,18,13]]_2$     &       $[[76,18,12]]_2$        \\
			16  &     $[[40,6,9]]_2$      &        $[[40,6,8]]_5$         & 75  &    $[[76,17,13]]_2$     &       $[[76,17,12]]_2$        \\
			17  &     $[[41,6,9]]_2$      &        $[[41,6,8]]_6$         & 76  &    $[[76,16,13]]_2$     &       $[[76,16,12]]_2$        \\
			18  &    $[[42,6,10]]_2^*$    &        $[[42,6,9]]_2$         & 77  &    $[[76,15,13]]_2$     &       $[[76,15,12]]_2$        \\
			19  &     $[[43,6,10]]_2$     &        $[[43,6,9]]_2$         & 78  &    $[[77,18,13]]_2$     &       $[[77,18,12]]_2$        \\
			20  &     $[[44,6,10]]_2$     &        $[[44,6,9]]_2$         & 79  &    $[[77,17,13]]_2$     &       $[[77,17,12]]_2$        \\
			21  &     $[[45,6,10]]_2$     &        $[[45,6,9]]_2$         & 80  &    $[[77,16,13]]_2$     &       $[[77,16,12]]_2$        \\
			22  &     $[[41,7,9]]_2$      &        $[[41,7,8]]_2$         & 81  &    $[[77,15,13]]_2$     &       $[[77,15,12]]_2$        \\
			23  &     $[[41,6,9]]_2$      &        $[[41,6,8]]_2$         & 82  &    $[[78,18,13]]_2$     &       $[[78,18,12]]_2$        \\
			24  &    $[[42,13,8]]_2^*$    &        $[[42,13,7]]_2$        & 83  &    $[[78,17,13]]_2$     &       $[[78,17,12]]_2$        \\
			25  &    $[[42,14,8]]_2^*$    &        $[[42,14,7]]_2$        & 84  &    $[[78,16,13]]_2$     &       $[[78,16,12]]_2$        \\
			26  &    $[[45,4,11]]_2^*$    &        $[[45,4,10]]_2$        & 85  &    $[[78,15,13]]_2$     &       $[[78,15,12]]_2$        \\
			27  &     $[[46,4,11]]_2$     &        $[[46,4,10]]_2$        & 86  &    $[[79,18,13]]_2$     &       $[[79,18,12]]_2$        \\
			28  &     $[[47,4,11]]_2$     &        $[[47,4,10]]_2$        & 87  &    $[[79,17,13]]_2$     &       $[[79,17,12]]_2$        \\
			29  &    $[[45,5,11]]_2^*$    &        $[[45,5,10]]_2$        & 88  &    $[[79,16,13]]_2$     &       $[[79,16,12]]_2$        \\
			30  &     $[[46,5,11]]_2$     &        $[[46,5,10]]_2$        & 89  &    $[[80,18,13]]_2$     &       $[[80,18,12]]_2$        \\
			31  &     $[[47,5,11]]_2$     &        $[[47,5,10]]_2$        & 90  &    $[[80,17,13]]_2$     &       $[[80,17,12]]_2$        \\
			32  &     $[[48,5,11]]_2$     &        $[[48,5,10]]_2$        & 91  &    $[[80,16,13]]_2$     &       $[[80,16,12]]_2$        \\
			33  &    $[[45,6,10]]_2^*$    &        $[[45,6,9]]_2$         & 92  &    $[[81,18,13]]_2$     &       $[[81,18,12]]_2$        \\
			34  &    $[[45,10,9]]_2^*$    &        $[[45,10,8]]_2$        & 93  &    $[[72,19,12]]_2$     &       $[[72,19,11]]_2$        \\
			35  &    $[[48,3,12]]_2^*$    &        $[[48,3,11]]_2$        & 94  &    $[[72,18,12]]_2$     &       $[[72,18,11]]_2$        \\
			36  & $[[45,21,7]]_2^{\star}$ &        $[[45,21,6]]_2$        & 95  &   $[[78,25,11]]_2^*$    &       $[[78,25,10]]_2$        \\
			37  &    $[[48,5,11]]_2^*$    &        $[[48,5,10]]_2$        & 96  &   $[[78,25,12]]_2^*$    &       $[[78,25,10]]_2$        \\
			38  &    $[[49,4,12]]_2^*$    &        $[[49,4,11]]_2$        & 97  &    $[[78,24,12]]_2$     &       $[[78,24,11]]_2$        \\
			39  &    $[[51,8,11]]_2^*$    &        $[[51,8,10]]_2$        & 98  &    $[[78,23,12]]_2$     &       $[[78,23,11]]_2$        \\
			40  &     $[[51,7,11]]_2$     &        $[[51,7,10]]_2$        & 99  &    $[[78,22,12]]_2$     &       $[[78,22,11]]_2$        \\
			41  &     $[[51,6,11]]_2$     &        $[[51,6,10]]_2$        & 100 &    $[[78,21,12]]_2$     &       $[[78,21,11]]_2$        \\
			42  &     $[[50,9,10]]_2$     &        $[[50,9,9]]_2$         & 101 &    $[[79,25,12]]_2$     &       $[[79,25,11]]_2$        \\
			43  &    $[[51,9,11]]_2^*$    &        $[[51,9,10]]_2$        & 102 &    $[[79,24,12]]_2$     &       $[[79,24,11]]_2$        \\
			44  &    $[[50,10,10]]_2$     &        $[[50,10,9]]_2$        & 103 &    $[[79,23,12]]_2$     &       $[[79,23,11]]_2$        \\
			45  & $[[55,20,9]]_2^{\star}$ &               -               & 104 &    $[[79,22,12]]_2$     &       $[[79,22,11]]_2$        \\
			46  &   $[[56,11,11]]_2^*$    &       $[[56,11,10]]_2$        & 105 &    $[[80,25,12]]_2$     &       $[[80,25,11]]_2$        \\
			47  &    $[[57,11,11]]_2$     &       $[[57,11,10]]_2$        & 106 &    $[[80,24,12]]_2$     &       $[[80,24,11]]_2$        \\
			48  &    $[[58,11,11]]_2$     &       $[[58,11,10]]_2$        & 107 &    $[[80,23,12]]_2$     &       $[[80,23,11]]_2$        \\
			49  &    $[[60,6,13]]_2^*$    &        $[[60,6,12]]_2$        & 108 &    $[[81,25,12]]_2$     &       $[[81,25,11]]_2$        \\
			50  &   $[[60,10,12]]_2^*$    &       $[[60,10,11]]_2$        & 109 &    $[[81,24,12]]_2$     &       $[[81,24,11]]_2$        \\
			51  &     $[[60,9,12]]_2$     &        $[[60,9,11]]_2$        & 110 &    $[[82,25,12]]_2$     &       $[[82,25,11]]_2$        \\
			52  &    $[[61,10,12]]_2$     &       $[[61,10,11]]_2$        & 111 &    $[[77,26,11]]_2$     &       $[[77,26,10]]_2$        \\
			53  &   $[[60,11,12]]_2^*$    &       $[[60,11,11]]_2$        & 112 &    $[[77,25,11]]_2$     &       $[[77,25,10]]_2$        \\
			54  &    $[[61,11,12]]_2$     &       $[[61,11,11]]_2$        & 113 &    $[[77,24,11]]_2$     &       $[[77,24,10]]_2$        \\
			55  &    $[[59,12,11]]_2$     &       $[[59,12,10]]_2$        & 114 &   $[[78,26,11]]_2^*$    &       $[[78,26,10]]_2$        \\
			56  &    $[[62,6,14]]_2^*$    &        $[[62,6,13]]_2$        & 115 &    $[[79,26,11]]_2$     &       $[[79,26,10]]_2$        \\
			57  &     $[[62,5,14]]_2$     &        $[[62,5,13]]_2$        & 116 &   $[[78,29,10]]_2^*$    &        $[[78,29,9]]_2$        \\
			58  &     $[[61,7,13]]_2$     &        $[[61,7,12]]_2$        & 117 &   $[[105,80,6]]_2^*$    &       $[[105,80,5]]_2$        \\
			59  &     $[[61,6,13]]_2$     &        $[[61,6,12]]_2$        &  -  &            -            &               -               \\ \hline
		\end{tabular}
		\begin{tablenotes}    %这行要添加， 从这开始
			\footnotesize               %这行要添加
			\item[1] The QECCs with superscripts ``$*$" and ``$\star$" can be found in Table \ref{F_QECCs} and Table \ref{T_QECCs} as construction polynomials, respectively. The rest of the QECCs are derived by Lemma \ref{propagation_rules}.          %这行要添加
			\item[2] Some additions: 
			$[[55,20,9]]_2$ and $[[63,31,8]]_2$ were previously presented in our first version \cite{guan2021new}, but were not formally published in \cite{guan2022new}; here we formally submit them for publication.
			After completing the first version of this paper \cite{guan2022symplectic}, we note that in the very near time, Dastbasteh et al. \cite{dastbasteh2023polynomial} also obtained $[[45,6,10]]_2$, $[[45,10,9]]_2$ and $[[51,8,11]]_2$, and Seneviratne \cite{Grassltable} obtained $[[40,6,10]]_2$, respectively.
%			After completing the first version \cite{guan2022symplectic} of this paper, we notice that later \cite{dastbasteh2023polynomial} also gets $[[45,6,10]]_2$, $[[45,10,9]]_2$ and $[[51,8,11]]_2$.
%			   In addition, during our revision of this paper \cite{guan2022symplectic}, we noticed that Padmapani Seneviratne also obtained $[[40,6,10]]_2$ and updated it in the code table \cite{Grassltable}.        %这行要添加
		\end{tablenotes}            %这行要添加
	\end{threeparttable}       %这行要添加，到这里结束
\end{table*}



\begin{table*}[!t]
	\caption{One-generator symplectic self-orthogonal QC codes with $\ell=2$ and record-breaking  QECCs}\label{F_QECCs}
	\centering
	\footnotesize
	\begin{tabular}{ccccc}
		\hline
		No. &Codes      &Symplectic Dual   &                                                                                                                                                                                                                             $g(x),f_0(x),f_1(x)$                                                                                                                                                                                                                              &    Record-breaking  QECCs     \\ \hline
		 1  &  $[80,35]_2^s$  &  $[80,45,10]_2^s$   &                                                                                                                     \makecell[c]{$1^{2}0^{2}1^2$, $0^{4}1^{3}01^{3}0101^{3}0^{2}10^{2}1^{3}0101^{3}01^3$,\\
			$0101^{2}0^{3}1^{2}0^{5}10^{2}1010^{2}10^{5}1^{2}0^{3}1^{2}01$}                                                                                                                      & $[[40,5,10]]_2$  \\
		 2  &  $[84,36]_2^s$  &  $[84,48,10]_2^s$   &                                                                                                                                                \makecell[c]{$101^2,1^{3}0^{3}1^{2}01^{3}0101^{2}0^{3}1^{3}0^{3}1^{2}0101^{3}01^{2}0^{3}1^2$,\\$010^{2}10^{2}1^{2}0^{3}10101^{2}0^{2}1010^{2}1^{2}01010^{3}1^{2}0^{2}10^{2}1$}                                                                                                                                                 & $[[42,6,10]]_2$  \\
		 3  &  $[84,29]_2^s$  &   $[84,55,8]_2^s$   &                                                                                                                                                     \makecell[c]{$101010^{3}10^{3}1,1^{2}0^{2}101^{2}0^{2}1^{3}01^{2}01^{2}01^{3}01^{2}01^{2}01^{3}0^{2}1^{2}010^{2}1$,\\$0^{7}10^{3}1^{2}0^{4}10^{7}10^{4}1^{2}0^{3}1$}                                                                                                                                                      & $[[42,13,8]]_2$  \\
		 4  &  $[84,28]_2^s$  &   $[84,56,8]_2^s$   &                                                                                                                                                         \makecell[c]{$10^{3}10^{3}10101,0^{4}10^{3}10^{3}101^{3}0^{2}10^{3}10^{2}1^{3}010^{3}10^{3}1$,\\$01^{3}0101^{6}0^{3}10^{2}1^{2}01^{2}0^{2}10^{3}1^{6}0101^3$}                                                                                                                                                         & $[[42,14,8]]_2$  \\
		 5  &  $[90,41]_2^s$  &  $[90,49,11]_2^s$   &                                                                                                                                                        \makecell[c]{$10^{2}1^2,1010101^{3}0^{2}10^{2}1^{3}01^{10}01^{3}0^{2}10^{2}1^{3}0101$,\\$101^{4}0^{3}1^{2}01010^{2}101^{3}0^{2}1^{3}010^{2}10101^{2}0^{3}1^4$ }                                                                                                                                                        & $[[45,4,11]]_2$  \\
		 6  &  $[90,40]_2^s$  &  $[90,50,11]_2^s$   &                                                                                                                                                  \makecell[c]{$10101^2,1^{2}0^{5}1010101^{2}0^{2}1^{2}0^{8}1^{2}0^{2}1^{2}0101010^{5}1$,\\$1^{2}0^{2}1^{3}0^{2}101^{3}010^{3}1^{3}0^{2}1^{3}0^{3}101^{3}010^{2}1^{3}0^{2}1$}                                                                                                                                                  & $[[45,5,11]]_2$  \\
		 7  &  $[90,39]_2^s$  &  $[90,51,10]_2^s$   &                                                                                                                                                   \makecell[c]{$10^{2}10^{2}1,10101^{2}0^{3}10^{2}10^{3}10^{2}10^{6}10^{2}10^{3}10^{2}10^{3}1^{2}01$,\\$1^{2}01^{3}0^{2}1010^{6}1010^{2}1^{2}0^{2}1010^{6}1010^{2}1^{3}01$}                                                                                                                                                   & $[[45,6,10]]_2$  \\
		 8  &  $[90,35]_2^s$  &   $[90,55,9]_2^s$   &                                                                                                                                                                                $1^5,1^{2}01^{2}010^{3}1^{4}010101^{3}0^{2}1^{3}010101^{4}0^{3}101^{2}01,10^{5}10^{5}10^{5}1^{10}0^{5}10^{5}1$                                                                                                                                                                                 & $[[45,10,9]]_2$  \\
		 9  &  $[96,45]_2^s$  &  $[96,51,12]_2^s$   &                                                                                                                                                      \makecell[c]{$1^2,010^{2}10^{9}1^{2}01010^{2}1^{2}01^{2}0^{2}10101^{2}0^{9}10^{2}1$,\\$01^{2}010^{2}101^{5}0^{2}1^{3}01^{2}0^{5}1^{2}01^{3}0^{2}1^{5}010^{2}101^2$}                                                                                                                                                      & $[[48,3,12]]_2$  \\
		 10  &  $[96,43]_2^s$  &  $[96, 53,11]_2^s$  &                                                                                                                                                       \makecell[c]{$1^2,0^{4}1^{3}0101^{3}01^{7}01^{2}01^{2}01^{7}01^{3}0101^3$,\\$0^{4}10^{2}10^{2}1^{3}01^{2}01010^{2}1^{2}01^{2}0^{2}10101^{2}01^{3}0^{2}10^{2}1$  }                                                                                                                                                       & $[[48,5,11]]_2$  \\
		11  & $[98, 45]_2^s$  &  $[98, 53,12]_2^s$  &                                                                                                                                       \makecell[c]{$101^2,010^{2}1^{2}0^{2}1^{2}010^{2}10^{6}1^{2}0^{4}1^{2}0^{6}10^{2}101^{2}0^{2}1^{2}0^{2}1$,\\$0^{3}1^{2}0^{4}10^{2}101^{2}01010^{2}1^{2}0^{2}1^{2}0^{2}10101^{2}010^{2}10^{4}1^2$}                                                                                                                                       & $[[49,4,12]]_2$  \\
		12  & $[102, 43]_2^s$ & $[102, 59,11]_2^s$  &                                                                                                                                     \makecell[c]{$10^{3}1^{2}01^2,1^{2}0101^{3}01^{2}0^{4}10^{3}1^{2}0101^{4}0101^{2}0^{3}10^{4}1^{2}01^{3}0101$,\\$1^{3}01^{6}0^{2}1^{2}010^{3}1^{2}0^{2}101^{2}010^{2}1^{2}0^{3}101^{2}0^{2}1^{6}01^2$}                                                                                                                                     & $[[51,8,11]]_2$  \\
		13  & $[102, 42]_2^s$ & $[102, 60,11]_2^s$  &                                                                                                                                            \makecell[c]{$10^{3}1^{2}01^2,0^{2}1010^{3}1^{7}01^{2}01^{3}0^{2}10^{2}10^{2}1^{3}01^{2}01^{7}0^{3}101$,\\$01^{5}0^{2}10^{4}1^{2}0^{4}10101^{2}0^{2}1^{2}01010^{4}1^{2}0^{4}10^{2}1^5$}                                                                                                                                            & $[[51,9,11]]_2$  \\
		14  & $[112, 45]_2^s$ & $[112, 67,11]_2^s$  &                                                                                                                                   \makecell[c]{$1^{2}0^{2}1^6,101010^{3}101^{2}0^{3}1^{4}0101^{2}0^{3}1^{3}0^{3}1^{2}0101^{4}0^{3}1^{2}010^{3}101$,\\$1^{4}0^{2}10^{9}1^{2}0^{2}1^{2}01^{3}0^{2}10^{2}1^{3}01^{2}0^{2}1^{2}0^{9}10^{2}1^3$}                                                                                                                                   & $[[56,11,11]]_2$ \\
		15  & $[120, 54]_2^s$ & $[120, 66,13]_2^s$  &                                                                                                                                         \makecell[c]{$1^{2}0^{2}1,0101010^{3}1^{5}01^{3}01^{2}0^{2}10^{3}1^{2}0^{3}1^{2}0^{3}10^{2}1^{2}01^{3}01^{5}0^{3}10101$,\\$1^{5}01^{4}0^{5}1010^{2}1^{5}01^{9}01^{5}0^{2}1010^{5}1^{4}01^4$}                                                                                                                                          & $[[60,6,13]]_2$  \\
		16  & $[120, 50]_2^s$ & $[120, 70,12]_2^s$  &                                                                                                                         \makecell[c]{$10^{5}101,0^{2}10^{3}10^{2}1^{2}01^{4}010^{4}1^{2}0^{2}1^{2}0^{5}1^{2}0^{2}1^{2}0^{4}101^{4}01^{2}0^{2}10^{3}1$,\\$01^{3}01^{3}01^{2}0^{3}101^{3}010^{2}1^{3}010^{5}101^{3}0^{2}101^{3}010^{3}1^{2}01^{3}01^3$}                                                                                                                         & $[[60,10,12]]_2$ \\
		17  & $[120, 49]_2^s$ & $[120, 71,12]_2^s$  &                                                                                                                           \makecell[c]{$10^{5}101,1^{2}010^{3}1^{4}0^{2}1010^{2}1010^{3}101^{3}0101^{3}010^{3}1010^{2}1010^{2}1^{4}0^{3}101$,\\$0^{3}10^{3}1^{3}01^{3}0^{3}1^{2}0^{2}10^{3}1^{2}0^{7}1^{2}0^{3}10^{2}1^{2}0^{3}1^{3}01^{3}0^{3}1$}                                                                                                                            & $[[60,11,12]]_2$ \\
		18  & $[124, 56]_2^s$ & $[124, 68,14]_2^s$  &                                                                                                                        \makecell[c]{$1^{4}01,1^{2}0^{2}101^{4}0^{7}1^{3}0^{4}10^{2}10^{3}10^{3}10^{2}10^{4}1^{3}0^{7}1^{4}010^{2}1$,\\$10^{2}1^{2}0^{3}1^{2}01^{4}010^{2}1010^{2}1^{2}01^{2}0^{2}10^{2}1^{2}01^{2}0^{2}1010^{2}101^{4}01^{2}0^{3}1^2$}                                                                                                                        & $[[62,6,14]]_2$  \\
		19  & $[146, 55]_2^s$ & $[146, 91,13]_2^s$  &                                                                                                     \makecell[c]{$1^{2}0^{3}10^{2}1^{3}0^{4}1^{2}01$,\\$10101^{3}0^{7}1^{2}010^{3}10^{2}1^{2}0^{2}1^{6}010^{2}101^{6}0^{2}1^{2}0^{2}10^{3}101^{2}0^{7}1^{3}01$,\\$1^{6}0^{2}101^{3}0^{2}10^{2}10^{2}1^{2}0^{6}1^{4}0^{8}1^{4}0^{6}1^{2}0^{2}10^{2}10^{2}1^{3}010^{2}1^5$}                                                                                                     & $[[73,18,13]]_2$ \\
		20  & $[156, 53]_2^s$ & $[156, 103,11]_2^s$ &                                                                                         \makecell[c]{$10101010^{7}10101010^{3}1$,\\$101^{2}0^{3}1^{2}0^{5}10^{3}1010^{2}10^{3}1^{3}010^{2}10^{4}10^{4}10^{2}101^{3}0^{3}10^{2}1010^{3}10^{5}1^{2}0^{3}1^2$,\\$01010^{2}101^{2}0^{5}10^{2}101^{4}0^{4}10^{2}101^{3}01^{2}01^{2}01^{3}010^{2}10^{4}1^{4}010^{2}10^{5}1^{2}010^{2}101$}                                                                                          & $[[78,25,11]]_2$ \\
		21  & $[156, 53]_2^s$ & $[156, 103,12]_2^s$ &                                                                                                 \makecell[c]{$10101010^{7}10101010^{3}1$,\\$10^{2}101^{2}01^{6}0101^{5}0^{2}1^{2}01^{2}01^{7}0^{2}10^{2}1^{7}01^{2}01^{2}0^{2}1^{5}0101^{6}01^{2}01$,\\$1^{5}0^{2}10^{3}1010^{4}1^{2}0^{2}1^{4}0^{5}10^{2}1^{2}01^{5}01^{2}0^{2}10^{5}1^{4}0^{2}1^{2}0^{4}1010^{3}10^{2}1^4$}                                                                                                 & $[[78,25,12]]_2$ \\
		22  & $[156, 52]_2^s$ & $[156, 104,11]_2^s$ &                                                                                                     \makecell[c]{$10101010^{7}10101010^{3}1$,\\$0^{3}101^{3}0^{5}10^{4}1010^{2}10^{4}1^{2}01010^{2}1^{3}01^{3}0^{2}10101^{2}0^{4}10^{2}1010^{4}10^{5}1^{3}01$,\\$01^{9}01^{2}0^{3}10101^{2}0^{3}10^{2}10^{2}101^{2}0^{9}1^{2}010^{2}10^{2}10^{3}1^{2}01010^{3}1^{2}01^9$}                                                                                                     & $[[78,26,11]]_2$ \\
		23  & $[156, 49]_2^s$ & $[156, 107,10]_2^s$ &                                                                                \makecell[c]{$10101010^{7}10101010^{3}1$,\\$10^{3}10101^{2}0^{5}1^{2}010^{2}10^{4}1010^{2}1^{2}0^{5}1^{3}0^{5}1^{2}0^{2}1010^{4}10^{2}101^{2}0^{5}1^{2}0101$,\\ $0^{2}101^{2}010^{2}1^{2}0^{5}1^{3}01^{2}01^{2}01^{2}0^{2}10^{2}10^{2}10^{3}10^{2}10^{2}10^{2}1^{2}01^{2}01^{2}01^{3}0^{5}1^{2}0^{2}101^{2}01$}                                                                                & $[[78,29,10]]_2$ \\
		24  & $[210, 25]_2^s$ & $[210, 185,6]_2^s$  & \makecell[c]{$10^{2}1^{3}0^{2}1^{4}010^{5}1^{4}0^{2}1^{4}0^{3}10^{4}1^{4}01^{3}01^{2}01010^{5}1^{2}010^{4}101^{3}0^{5}1^{2}0^{2}1^2$,\\$10^{2}1^{2}0^{2}101^{5}0^{4}101^{2}0^{2}1^{2}01^{4}01^{3}0101^{3}0^{2}1^{2}0^{2}1^{2}0^{2}1^{2}0^{2}$\\$10^{2}10^{2}1^{4}010^{7}1^{2}0101^{3}01010^{5}1^{2}01^{2}0^{2}1010^{3}1,10^{2}1^{2}01^{2}0101^{2}01^{5}$\\$0^{3}1^{2}0^{2}10^{6}1^{4}0^{4}1010101^{3}0101^{3}0^{2}101^{2}010^{2}10^{7}1^{4}0^{4}1^{3}01^{4}0^{2}10^{3}101^2$} & $[[105,80,6]]_2$ \\ \hline
	\end{tabular}
\end{table*}


\begin{table*}[!t]
			\centering
		\begin{threeparttable}
	\caption{Two-generator symplectic self-orthogonal QC codes\tnote{3} with $\ell=2$  and record-breaking  QECCs}\label{T_QECCs}
	\centering
	\footnotesize
	\begin{tabular}{ccccc}	
		\hline
		No. &     Codes      &   Symplectic Dual  &                                                                                                                                $g_1(x),g_2(x),f(x)$                                                                                                                                 &    Record-breaking  QECCs    \\ \hline
		 1  & $[80,34]_2^s$  &  $[80,46,10]_2^s$   &                                                                                               \makecell[c]{$1^8$, $1^{40}$, $0^{4}1^{5}01^{2}0^{2}101^{3}0^{3}1^{3}010^{2}1^{2}01^5$}                                                                                               & $[[40,6,10]]_2$ \\
		 2  & $[84,35]_2^s$  &  $[84,49,10]_2^s$   &                                                                                               \makecell[c]{$10^{7}101$, $101^{3}0^{2}101^{3}0^{2}101^{3}0^{2}101^{3}0^{2}101^{3}0^{2}101^3$,\\
		 	$1^{3}0^{8}1^{3}0^{4}1^{7}0^{4}1^{3}0^{8}1^2$}                                                                                               & $[[42,7,10]]_2$ \\
		 3  & $[90,24]_2^s$  & $[ 90, 66, 7 ]_2^s$ &                                          \makecell[c]{$1^{4}0101^{4}0^{2}101^{2}0^{2}10^{3}1$,\\
	$1^{2}01^{2}01^{2}01^{2}01^{2}01^{2}01^{2}01
	^{2}01^{2}01^{2}01^{2}01^{2}01^{2}01^{2}01^2$,\\
	$0101^{2}0^{5}1010^{20}1010^{5}1^{2}01$}                                          & $[[45,21,7]]_2$ \\
		 4  & $[110,35]_2^s$ &  $[110,75,9]_2^s$   &                                                  \makecell[c]{$10^{2}101^{2}010^{2}10^{3}1^{2}0^{4}10^{2}10^{2}1010101^2$,\\
	$1010^{2}1010^{2}1^{4}01^{4}01^{2}01^{4}01^{4}0^{2}1010^{2}101$,\\
	$010^{3}1^{2}0^{42}1^{2}0^{3}1$}                                                  & $[[55,20,9]]_2$ \\
		 5  & $[126,32]_2^s$ &  $[126,94,8]_2^s$   & \makecell[c]{$10^{4}101^{2}0^{2}1^{2}0^{4}10^{2}10101010101^{2}01^{3}01^3$,\\
	$1010^{2}1^{2}0^{2}1^{2}010^{2}1^{3}0^{2}10^{2}1010^{3}1^{2}0101^{2}01010^{4}1^{3}0^{4}1^{2}0^{2}1$,\\ $0^{2}1^{2}0^{2}1^{3}010^{5}101^{2}010101^{2}010^{6}101^{2}010101^{2}010^{5}101^{3}0^{2}1^2$} & $[[63,31,8]]_2$ \\ \hline
	\end{tabular}
\begin{tablenotes}    %这行要添加， 从这开始
	\footnotesize               %这行要添加
	\item[3] These $2$-generator QC codes with generators $([g_1(x)f(x)],[g_1(x)])$ and $([g_2(x)],[g_2(x)f(x)])$.        %这行要添加
%	\item[4] $[[55,20,9]]_2$ and $[[63,31,8]]_2$ were presented in our first version of \cite{guan2021new}, but were not formally published in \cite{guan2022new}; here we formally submit them for publication.  In addition, during our revision of this paper \cite{guan2022symplectic}, we noticed that Padmapani Seneviratne also obtained $[[40,6,10]]_2$ and updated it in the code table \cite{Grassltable}. 
\end{tablenotes}            %这行要添加
\end{threeparttable}       %这行要添加，到这里结束
\end{table*}

 


%\section*{Acknowledgments}
%The authors would like to thank Professor Markus Grassl, Dr. Shitao Li, and Dr. Yang Li for their valuable suggestions, which greatly improved the presentation and quality of this paper.




%%%%%%%%%%%%%%%%%%%%%%%%%%%%%%%%%%%%%%%%%%%%%%%%%%%%%%%%%%%%%%%%%%%%%%%%%
\bibliographystyle{IEEEtran} 
\bibliography{reference} 
%%%%%%%%%%%%%%%%%%%%%%%%%%%%%%%%%%%%%%%%%%%%%%%%%%%%%%%%%%%%%%%%%%%%%%%%%



%	\begin{IEEEbiographynophoto}{Chaofeng Guan}
%	is currently pursuing the Ph.D. degree with the Henan Key Laboratory of Network Cryptography in Zhengzhou, China. His research interests include coding theory and cryptography.
%\end{IEEEbiographynophoto}
%
%\begin{IEEEbiographynophoto}{Ruihu Li}
%	received Ph.D. degree from Northwestern Polytechnical University of Technology in 2004. He is currently a professor in the Fundamentals Department, Air Force Engineering University, Xi'an, China. His research interests include group theory, coding theory and cryptography.
%\end{IEEEbiographynophoto}
%
%\begin{IEEEbiographynophoto}{Jingjie Lv}
%	received the B.S. and M.S. degrees in mathematics and applied mathematics from the Shandong University of Technology, Zibo, China, in 2015 and 2018, respectively, and the Ph.D. degree in cyberspace security from the Quantum Information Laboratory, Air Force Engineering University, Xi’an, China, in 2021. From 2021 to 2023, he was a PostDoctoral Researcher with the Tsinghua Shenzhen International Graduate School, Tsinghua University, Shenzhen, China. Since November 2023, he has been an Assistant Researcher with School of Electrical Engineering \& Intelligentization, Dongguan University of Technology, Dongguan, China. His research interests include quantum coding and classical coding for communications and distributed storage systems.
%\end{IEEEbiographynophoto}
%
%
%\begin{IEEEbiographynophoto}{Zhi Ma}
%	received the Ph.D. degree in Mathematics from the University of Science and Technology of China in 2002. She is currently a professor at the Henan Key Laboratory of Network Cryptography Technology, Zhengzhou, China. Her research interests include quantum information and quantum computation.
%\end{IEEEbiographynophoto}
\end{document}


 
